\documentclass[11pt]{article}
\usepackage[margin=1in]{geometry}
\usepackage[utf8]{inputenc}
\DeclareUnicodeCharacter{0233}{\ensuremath{\bar{y}}}
\DeclareUnicodeCharacter{221A}{\ensuremath{\surd}}
\DeclareUnicodeCharacter{03A3}{\ensuremath{\Sigma}}
\DeclareUnicodeCharacter{2192}{\ensuremath{\rightarrow}}
\DeclareUnicodeCharacter{2265}{\ensuremath{\geq}}
\DeclareUnicodeCharacter{2211}{\ensuremath{\sum}}
\DeclareUnicodeCharacter{2003}{ }
\DeclareUnicodeCharacter{2212}{\ensuremath{-}}
\DeclareUnicodeCharacter{03A6}{\ensuremath{\Phi}}
\DeclareUnicodeCharacter{2264}{\ensuremath{\leq}}
\DeclareUnicodeCharacter{2265}{\ensuremath{\geq}}
\DeclareUnicodeCharacter{00D7}{\ensuremath{\times}}
\DeclareUnicodeCharacter{2248}{\ensuremath{\approx}}
\DeclareUnicodeCharacter{03B2}{\ensuremath{\beta}}
\DeclareUnicodeCharacter{03C4}{\ensuremath{\tau}}
\DeclareUnicodeCharacter{03C7}{\ensuremath{\chi}}
\DeclareUnicodeCharacter{0302}{}
\DeclareUnicodeCharacter{0304}{}
\DeclareUnicodeCharacter{00B7}{\ensuremath{\cdot}}
\DeclareUnicodeCharacter{2032}{\ensuremath{'}}
\usepackage{textcomp}
\usepackage{graphicx,booktabs,longtable,array,calc,float,xcolor,setspace,ragged2e}\usepackage[expansion=false]{microtype}
\usepackage[skins]{tcolorbox}
\usepackage[font=small,labelformat=empty]{caption}
\usepackage[hidelinks]{hyperref}
\usepackage{titlesec}\usepackage{pdflscape}\graphicspath{{/tmp/claude-0/-home-claude/1e37972c-6107-52a0-86b9-b3cd81bfbeeb/scratchpad/jama/}}
\titleformat{\section}{\large\bfseries}{}{0pt}{}
\titleformat{\subsection}{\normalsize\bfseries}{}{0pt}{}
\titlespacing*{\section}{0pt}{1.4em}{0.6em}
\titlespacing*{\subsection}{0pt}{1.1em}{0.4em}
\setstretch{1.25}
\setlength{\parskip}{0.35em}
\providecommand{\tightlist}{\setlength{\itemsep}{0pt}\setlength{\parskip}{0pt}}
\newcommand{\passthrough}[1]{#1}
\def\real#1{#1}
\setlength{\LTleft}{0pt}\setlength{\LTright}{0pt}
\begin{document}
\begin{center}{\LARGE Supplemental Online Content\par}\vspace{1em}{\large Early-Life Infections and Later Autism and Attention-Deficit/Hyperactivity Disorder\par}\vspace{0.8em}Ssentongo AE, Ssentongo P.\par\vspace{0.4em}2026\end{center}\vspace{1em}\noindent\textit{This supplemental material has been provided by the authors to give readers additional information about their work.}

\section*{Supplement Contents}

\begin{itemize}
\tightlist
\item
  eMethods 1. Data Source and Privacy Protections
\item
  eMethods 2. Cohort Construction
\item
  eMethods 3. Covariates and Propensity Score Model
\item
  eMethods 4. Matching Algorithm
\item
  eMethods 5. Outcome Ascertainment and Time-to-Event
  Analysis
\item
  eMethods 6. Derived Statistics: P Values, False Discovery
  Rate, Absolute Risks, Calibration, Pooling, and E-Values
\item
  eMethods 7. Sensitivity Analyses and Deviations From the
  Analysis Plan
\item
  eTable 1. Exposure Definitions, Cohort Sizes, and Matched
  Pairs
\item
  eTable 2. Propensity Score Covariates and Codes
\item
  eTable 3. Outcome Definitions and Codes
\item
  eTable 4. Full Results for All Conditions and Outcomes
\item
  eTable 5. Hazard Ratios Calibrated Against Diaper
  Dermatitis
\item
  eTable 6. Random-Effects Pooled Hazard Ratios by Condition
  Group
\item
  eTable 7. Sensitivity Analysis Without the Well-Child
  Requirement in Exposed Infants
\item
  eTable 8. E-Values for Statistically Significant Positive
  Associations
\item
  eTable 9. Complete Baseline Characteristics for the CMV
  Comparison
\item
  eTable 10. Same-Era Sensitivity Analysis for COVID-19
\item
  eFigure 1. Calibrated Hazard Ratios for Autism and ADHD
\item
  eFigure 2. Absolute Risk Differences per 1,000 Children
\item
  eFigure 3. Primary vs Unrestricted Design
\item
  eFigure 4. E-Values
\item
  eReferences
\end{itemize}

\newpage

\hypertarget{emethods-1.-data-source-and-privacy-protections}{%
\section*{eMethods 1. Data Source and Privacy
Protections}\label{emethods-1.-data-source-and-privacy-protections}}

TriNetX is a federated health research network that links
deidentified electronic health record (EHR) data from
participating health care organizations (HCOs), mostly
academic medical centers and their affiliated hospitals and
clinics in the US.\textsuperscript{1} Data include
demographics, diagnoses coded in the \emph{International
Classification of Diseases, Tenth Revision, Clinical
Modification} (ICD-10-CM), procedures, medications mapped to
the Anatomical Therapeutic Chemical (ATC) classification and
RxNorm, and laboratory results. Queries are run in real time
on each HCO's data; only aggregate counts and model results
are returned to the investigator. TriNetX rounds any count
of 1 to 10 to 10, and it restricts analysis of selected
codes considered sensitive (here, congenital cytomegalovirus
{[}P35.1{]}, neonatal bacterial sepsis {[}P36{]}, and
neonatal jaundice {[}P59{]}). The analyses used the Global
Collaborative Network in study 409549; all analyses were run
in October 2026. The number of contributing HCOs varies by
query and over time. At the time of the final queries
(October 26, 2026), 188 of 193 network HCOs were online,
covering about 211 million patients; 150 HCOs contributed
infants to the well-child source population, 94\% of whom
were seen at US-based HCOs.

\hypertarget{emethods-2.-cohort-construction}{%
\section*{eMethods 2. Cohort
Construction}\label{emethods-2.-cohort-construction}}

\textbf{Source population.} Infants with at least 1 newborn,
infant, or child health examination (ICD-10-CM Z00.1)
recorded at age 0-1 year in completed years, that is, before
the second birthday (TriNetX age-at-event range 0-1).

\textbf{Exposure cohorts (22).} For each condition, the
cohort query required (1) at least 1 diagnosis code for the
condition (eTable 1) recorded at age 0-1 year (before age 2
years) and (2) Z00.1 recorded at age 0-1 year. The index
event was the first occurrence of the condition code only;
Z00.1 could occur before or after the index date within the
age window. Twenty-two cohort queries were built (W00-W23,
omitting the 2 restricted neonatal conditions).

\textbf{Control cohorts.} For each comparison, the control
query required Z00.1 at age 0-1 year and excluded any of
that comparison's exposure codes at age 0-1 year. The index
event was the Z00.1 visit. Because each control query
excludes only the codes of its own comparison, controls in
one comparison may carry other conditions (for example, a
control in the RSV comparison may have had influenza). This
mirrors the target question for each condition: risk
compared with infants in routine care who did not have that
condition.

\textbf{Why the well-child requirement.} Controls entered
through a routine visit and were therefore in routine care
by definition. Exposed infants identified only by a
sick-visit diagnosis include infants who received most of
their care elsewhere or only in acute settings. Such infants
have fewer recorded developmental assessments, which lowers
recorded rates of conditions such as ADHD and
speech-language disorder, and they include sicker infants
seen in hospitals, which raises recorded rates of severe
outcomes. Requiring Z00.1 in both groups aligns the
opportunity for detection. The unrestricted design is
reported as a sensitivity analysis (eTable 7; eFigure 3).

\hypertarget{emethods-3.-covariates-and-propensity-score-model}{%
\section*{eMethods 3. Covariates and Propensity Score
Model}\label{emethods-3.-covariates-and-propensity-score-model}}

Twenty-seven covariates were prespecified (eTable 2): 8
demographic variables (current age, age at index, male sex,
Hispanic or Latino ethnicity, not Hispanic or Latino
ethnicity, White race, Black or African American race, Asian
race), 11 perinatal and congenital diagnosis groups, 6
early-childhood diagnosis groups, 1 family-history code, and
1 medication class. Covariates were assessed in a window
from any time before through the index date. Diagnosis and
medication covariates were coded as binary indicators
(present at least once in the window vs not). Age variables
were continuous. Female sex and unknown or other race and
ethnicity categories were omitted as reference categories.

The original plan included 30 covariates. TriNetX limits the
computational size of the matching step (approximately the
product of the combined cohort size and the number of
covariates; about 100 million). Because several comparisons
exceeded this limit with 30 covariates, 3 covariates with
the weakest expected confounding role were removed for all
comparisons: feeding problems of newborn (P92),
gastroesophageal reflux disease (K21), and family history of
disabling chronic disease (Z82). Their balance after
matching is shown for the CMV comparison (Table 1 and eTable
9).

The propensity score for infant \emph{i} was estimated as

logit Pr(\emph{E\textsubscript{i}} = 1 \textbar{}
\textbf{\emph{X}}\emph{\textsubscript{i}}) =
β\textsubscript{0} +
β′\textbf{\emph{X}}\emph{\textsubscript{i}},

where \emph{E\textsubscript{i}} indicates membership in the
exposure cohort and
\textbf{\emph{X}}\emph{\textsubscript{i}} is the vector of
the 27 covariates, entered as main effects without
interactions or splines. The model was fitted by maximum
likelihood by the TriNetX platform (Python scikit-learn
logistic regression implementation).

\hypertarget{emethods-4.-matching-algorithm}{%
\section*{eMethods 4. Matching
Algorithm}\label{emethods-4.-matching-algorithm}}

Matching used 1:1 greedy nearest-neighbor matching without
replacement on the propensity score, with a caliper of 0.1
pooled SDs of the propensity score. TriNetX randomizes the
order of records before matching to avoid bias from the
order of rows in the data. An exposed infant was discarded
if no control fell within the caliper. Balance was
summarized with the absolute standardized mean difference
(SMD):

for a binary covariate, SMD =
\textbar{}\emph{p}\textsubscript{1} −
\emph{p}\textsubscript{2}\textbar{} /
√\{{[}\emph{p}\textsubscript{1}(1 −
\emph{p}\textsubscript{1}) + \emph{p}\textsubscript{2}(1 −
\emph{p}\textsubscript{2}){]} / 2\};

for a continuous covariate, SMD =
\textbar{}\emph{x̄}\textsubscript{1} −
\emph{x̄}\textsubscript{2}\textbar{} /
√{[}(\emph{s}\textsubscript{1}\textsuperscript{2} +
\emph{s}\textsubscript{2}\textsuperscript{2}) / 2{]}.

SMD values below 0.10 were taken to indicate adequate
balance. The SMDs in Table 1 and eTable 9 were recalculated
by the authors from the TriNetX baseline counts; cells
reported by TriNetX as 10 (true value 1-10) make the
corresponding SMDs approximate.

Matching was performed once per comparison. For each
outcome, infants with that outcome on or before the index
date were then excluded from both matched groups; matched
pairs were not re-formed, which is why at-risk numbers
differ slightly between groups and outcomes (eTable 4).

\hypertarget{emethods-5.-outcome-ascertainment-and-time-to-event-analysis}{%
\section*{eMethods 5. Outcome Ascertainment and Time-to-Event
Analysis}\label{emethods-5.-outcome-ascertainment-and-time-to-event-analysis}}

Outcomes were identified by the first ICD-10-CM code in the
follow-up window from 1 day to 3,285 days (9 years) after
the index date (eTable 3). Starting follow-up on day 1
excludes diagnoses recorded at the index encounter. Infants
were censored at their last recorded fact in the EHR or at
the end of the window, whichever came first.

Within each matched cohort, the hazard ratio (HR) was
estimated by a Cox proportional hazards model with exposure
group as the only covariate:

\emph{h}(\emph{t} \textbar{} \emph{E}) =
\emph{h}\textsubscript{0}(\emph{t}) exp(β\emph{E}), HR =
exp(β̂),

with 95\% CIs from the model-based standard error. Because
confounding was addressed by matching, no further covariates
were added, and matched pairs were not modeled as strata,
which yields a slightly conservative variance. Kaplan-Meier
survival functions were estimated for each group and
compared by the log-rank test. TriNetX also reports a test
of proportional hazards based on scaled Schoenfeld
residuals; these were not extracted for all 173 estimable
models, so HRs are interpreted as weighted averages of the
hazard ratio over follow-up.

\hypertarget{emethods-6.-derived-statistics}{%
\section*{eMethods 6. Derived
Statistics}\label{emethods-6.-derived-statistics}}

All derived statistics were computed by the authors in
Python 3.13 with SciPy from the TriNetX outputs (numbers at
risk, numbers of events, HR, and 95\% CI).

\textbf{Standard errors and P values.} For each HR with 95\%
CI (\emph{L}, \emph{U}), SE(ln HR) = (ln \emph{U} − ln
\emph{L})/(2 × 1.96), \emph{z} = ln HR / SE, and the 2-sided
\emph{P} = 2{[}1 − Φ(\textbar{}\emph{z}\textbar){]}, where Φ
is the standard normal distribution function.

\textbf{Absolute risks.} Cumulative risk per 1,000 children
= 1,000 × events / number at risk, computed for each group;
the risk difference (RD) is exposed minus control. These are
crude proportions over the follow-up window and do not
account for censoring; Kaplan-Meier estimates would be
somewhat higher.

\textbf{False discovery rate.} The Benjamini-Hochberg
procedure\textsuperscript{2} was applied to the family of
144 tests formed by 21 conditions (excluding the reference
exposure) × 7 outcomes (excluding the negative-control
outcome), less 3 tests that could not be estimated because
of suppressed counts. With ordered \emph{P} values
\emph{P}\textsubscript{(1)} ≤ \ldots{} ≤
\emph{P}\textsubscript{(m)}, \emph{q}\textsubscript{(i)} =
min\textsubscript{j≥i} \{\emph{m P}\textsubscript{(j)} /
\emph{j}\}, capped at 1.

\textbf{Calibration against the reference exposure.} For
condition \emph{c} and outcome \emph{o}, ln
HR\textsuperscript{cal}\textsubscript{c,o} = ln
HR\textsubscript{c,o} − ln HR\textsubscript{ref,o} and
SE\textsuperscript{cal} = √{[}SE(ln
HR\textsubscript{c,o})\textsuperscript{2} + SE(ln
HR\textsubscript{ref,o})\textsuperscript{2}{]}, where
\emph{ref} is diaper dermatitis. The 95\% CI is exp(ln
HR\textsuperscript{cal} ± 1.96 SE\textsuperscript{cal}). The
variance treats the 2 estimates as independent. Because the
2 comparisons share part of their control pool, the true
covariance is positive, so the true variance is smaller and
the reported CIs are conservative.

\textbf{Random-effects pooling.} Within each condition group
with at least 2 conditions, log HRs
\emph{y\textsubscript{j}} with within-study variances
\emph{v\textsubscript{j}} were pooled by the
DerSimonian-Laird method: fixed-effect weights
\emph{w\textsubscript{j}} = 1/\emph{v\textsubscript{j}},
Cochran \emph{Q} =
Σ\emph{w\textsubscript{j}}(\emph{y\textsubscript{j}} −
\emph{ȳ}\textsubscript{FE})\textsuperscript{2},
τ\textsuperscript{2} = max\{0, (\emph{Q} −
\emph{df})/{[}Σ\emph{w\textsubscript{j}} −
Σ\emph{w\textsubscript{j}}\textsuperscript{2}/Σ\emph{w\textsubscript{j}}{]}\},
random-effects weights \emph{w\textsubscript{j}}* =
1/(\emph{v\textsubscript{j}} + τ\textsuperscript{2}), and
pooled ln HR =
Σ\emph{w\textsubscript{j}}*\emph{y\textsubscript{j}} /
Σ\emph{w\textsubscript{j}}* with SE =
(Σ\emph{w\textsubscript{j}}*)\textsuperscript{−1/2}.
Heterogeneity was summarized by \emph{I}\textsuperscript{2}
= max\{0, (\emph{Q} − \emph{df})/\emph{Q}\}. Pooled
estimates describe the average association within a group of
related conditions; they do not imply that the conditions
are interchangeable, and high \emph{I}\textsuperscript{2}
values (eTable 6) indicate that individual estimates (eTable
4) should be preferred.

\textbf{E-values.} For an HR treated as an approximate risk
ratio \emph{RR} (reasonable because the outcomes are
uncommon), E-value = \emph{RR} + √{[}\emph{RR}(\emph{RR} −
1){]} for \emph{RR} ≥ 1 (with \emph{RR} replaced by
1/\emph{RR} when \emph{RR} \textless{}
1).\textsuperscript{3} The E-value for the CI uses the limit
closest to 1 and equals 1 if the CI includes 1.

\hypertarget{emethods-7.-sensitivity-analyses-and-deviations-from-the-analysis-plan}{%
\section*{eMethods 7. Sensitivity Analyses and Deviations
From the Analysis
Plan}\label{emethods-7.-sensitivity-analyses-and-deviations-from-the-analysis-plan}}

\begin{enumerate}
\def\labelenumi{\arabic{enumi}.}
\tightlist
\item
  \textbf{Unrestricted design.} Eighteen comparisons were
  also run without the Z00.1 requirement in the exposure
  cohort (controls unchanged). Bacterial meningitis, viral
  meningitis or encephalitis, and HSV infection used 30
  covariates in this analysis; the others used the 27
  covariates. Unrestricted analyses for febrile seizure,
  atopic dermatitis, and iron-deficiency anemia were not
  completed.
\item
  \textbf{Earlier CMV analysis.} A separate CMV analysis
  (B25 at age 0-1 year vs well-child controls without B25)
  was matched on 9 covariates (age at index, current age,
  sex, race, ethnicity, P07, Q20-Q28, D80-D89, Q90-Q99),
  yielding 5,147 pairs.
\item
  \textbf{Restricted codes and congenital CMV overlap.}
  Congenital CMV (P35.1), neonatal sepsis (P36), and
  neonatal jaundice (P59) were planned exposures but cannot
  define analysis cohorts in TriNetX. P35.1 can, however, be
  counted. Of 1,387 infants meeting the CMV exposure
  definition (B25 and Z00.1 before age 2 years), 635
  (45.8\%) also carried P35.1 (60 HCOs). Conversely, 1,492
  infants had P35.1 and a well-child visit before age 2
  years, of whom 635 (42.6\%) also had B25 and were
  therefore in the CMV cohort. The CMV cohort is thus a
  mixture of coded congenital CMV and CMV disease without a
  congenital code (postnatal infection or uncoded congenital
  infection).
\item
  \textbf{Negative-control outcome for bone infection.} Arm
  fracture is not a valid negative control for osteomyelitis
  or septic arthritis, because bone infection is linked to
  imaging, bone pathology, and coding of fractures.
\item
  \textbf{Same-era (calendar-time) analysis.} COVID-19
  (U07.1) can occur only from 2020 onward, whereas controls
  in the primary analysis were indexed across the full data
  period. The COVID-19 comparison was therefore repeated
  with both the exposure and control queries restricted to
  children with a current age of 0-6 years (born from late
  2019 onward), so that all well-child visits and index
  dates before age 2 years fell in the pandemic era. Because
  TriNetX does not offer an absolute calendar-date filter on
  cohort criteria, the current-age filter was used as the
  era restriction. The diaper dermatitis comparison was
  repeated with the same restriction to provide a same-era
  reference for calibration. All other settings (27
  covariates, matching, outcomes, and window) were
  unchanged. In both runs, TriNetX reported that 1 or more
  HCOs did not return results. Results are in eTable 10.
\end{enumerate}

\begin{landscape}\scriptsize

\hypertarget{etable-1.-exposure-definitions-cohort-sizes-and-matched-pairs}{%
\section*{eTable 1. Exposure Definitions, Cohort Sizes, and
Matched
Pairs}\label{etable-1.-exposure-definitions-cohort-sizes-and-matched-pairs}}

\begin{longtable}[]{@{}
  >{\raggedright\arraybackslash}p{(\columnwidth - 12\tabcolsep) * \real{0.2672}}
  >{\raggedright\arraybackslash}p{(\columnwidth - 12\tabcolsep) * \real{0.2672}}
  >{\raggedright\arraybackslash}p{(\columnwidth - 12\tabcolsep) * \real{0.1724}}
  >{\raggedleft\arraybackslash}p{(\columnwidth - 12\tabcolsep) * \real{0.0690}}
  >{\raggedleft\arraybackslash}p{(\columnwidth - 12\tabcolsep) * \real{0.0690}}
  >{\raggedleft\arraybackslash}p{(\columnwidth - 12\tabcolsep) * \real{0.0862}}
  >{\raggedleft\arraybackslash}p{(\columnwidth - 12\tabcolsep) * \real{0.0690}}@{}}
\toprule\noalign{}
\begin{minipage}[b]{\linewidth}\raggedright
Group
\end{minipage} & \begin{minipage}[b]{\linewidth}\raggedright
Condition
\end{minipage} & \begin{minipage}[b]{\linewidth}\raggedright
ICD-10-CM codes
\end{minipage} & \begin{minipage}[b]{\linewidth}\raggedleft
Infants with code before age 1 y, No.~\textsuperscript{a}
\end{minipage} & \begin{minipage}[b]{\linewidth}\raggedleft
With well-child visit, No.~\textsuperscript{b}
\end{minipage} & \begin{minipage}[b]{\linewidth}\raggedleft
Eligible controls, No.
\end{minipage} & \begin{minipage}[b]{\linewidth}\raggedleft
Matched pairs, No.
\end{minipage} \\
\midrule\noalign{}
\endhead
\bottomrule\noalign{}
\endlastfoot
Congenital/perinatal viral & CMV disease & B25 & 5,518 &
1,387 & 3,228,691 & 1,345 \\
Invasive CNS & Bacterial meningitis & G00 & 8,738 & 1,348 &
3,242,202 & 1,253 \\
Invasive CNS & Viral meningitis/encephalitis & A87, A86 &
11,166 & 1,859 & 3,241,670 & 1,685 \\
Invasive systemic & Sepsis (postneonatal) & A40, A41 &
68,893 & 10,654 & 3,232,794 & 9,667 \\
Invasive systemic & Osteomyelitis/septic arthritis & M86,
M00 & 9,492 & 1,381 & 3,242,163 & 1,262 \\
Invasive systemic & UTI/pyelonephritis & N39.0, N10 &
291,815 & 60,499 & 3,182,668 & 57,349 \\
Respiratory & RSV infection & J21.0, J12.1, B97.4 & 534,279
& 166,721 & 3,076,829 & 149,865 \\
Respiratory & Influenza & J09-J11 & 364,318 & 142,086 &
3,100,869 & 127,676 \\
Respiratory & COVID-19 & U07.1 & 254,511 & 95,708 &
3,147,123 & 94,108 \\
Respiratory & Pneumonia & J13, J15, J18 & 464,136 & 119,589
& 3,123,318 & 101,444 \\
Respiratory & Pertussis & A37 & 16,062 & 3,354 & 3,240,167 &
3,033 \\
Mucocutaneous/enteric/exanthem & HSV infection & B00 &
44,754 & 13,795 & 3,229,641 & 12,702 \\
Mucocutaneous/enteric/exanthem & Enterovirus infection &
B34.1, B97.1 & 135,422 & 40,698 & 3,201,085 & 39,922 \\
Mucocutaneous/enteric/exanthem & Viral gastroenteritis & A08
& 312,768 & 124,821 & 3,117,695 & 112,443 \\
Mucocutaneous/enteric/exanthem & Roseola & B08.2 & 45,130 &
12,410 & 3,231,125 & 12,265 \\
Mucocutaneous/enteric/exanthem & Varicella & B01 & 29,344 &
5,954 & 3,237,577 & 4,755 \\
Mucocutaneous/enteric/exanthem & Hand, foot, and mouth
disease & B08.4 & 295,982 & 154,545 & 3,085,392 & 132,566 \\
Noninfectious/inflammatory & Kawasaki disease & M30.3 &
16,970 & 1,881 & 3,241,658 & 1,825 \\
Noninfectious/inflammatory & Febrile seizure & R56.0 &
178,928 & 32,660 & 3,210,890 & 31,563 \\
Noninfectious/inflammatory & Atopic dermatitis & L20 &
693,145 & 396,999 & 2,846,557 & 367,367 \\
Noninfectious/inflammatory & Iron-deficiency anemia & D50 &
134,714 & 66,097 & 3,177,453 & 62,527 \\
Benign reference & Diaper dermatitis (reference) & L22 &
822,224 & 502,903 & 2,740,656 & 409,413 \\
\end{longtable}

\textsuperscript{a} First code before age 2 years, without
the well-child requirement (from the cohort query at the
time of analysis). \textsuperscript{b} Primary exposure
cohort: code and a well-child visit (Z00.1) before age 2
years. Planned but not analyzable because TriNetX restricts
the codes: congenital CMV (P35.1), neonatal sepsis (P36),
and neonatal jaundice (P59).

\newpage

\hypertarget{etable-2.-propensity-score-covariates-and-codes}{%
\section*{eTable 2. Propensity Score Covariates and
Codes}\label{etable-2.-propensity-score-covariates-and-codes}}

\begin{longtable}[]{@{}
  >{\raggedleft\arraybackslash}p{(\columnwidth - 6\tabcolsep) * \real{0.1484}}
  >{\raggedright\arraybackslash}p{(\columnwidth - 6\tabcolsep) * \real{0.2422}}
  >{\raggedright\arraybackslash}p{(\columnwidth - 6\tabcolsep) * \real{0.3047}}
  >{\raggedright\arraybackslash}p{(\columnwidth - 6\tabcolsep) * \real{0.3047}}@{}}
\toprule\noalign{}
\begin{minipage}[b]{\linewidth}\raggedleft
No.
\end{minipage} & \begin{minipage}[b]{\linewidth}\raggedright
Domain
\end{minipage} & \begin{minipage}[b]{\linewidth}\raggedright
Covariate
\end{minipage} & \begin{minipage}[b]{\linewidth}\raggedright
Code (ICD-10-CM, HL7/CDC race and ethnicity, or ATC)
\end{minipage} \\
\midrule\noalign{}
\endhead
\bottomrule\noalign{}
\endlastfoot
1 & Demographic & Current age (y), continuous & NA \\
2 & Demographic & Age at index (y), continuous & NA \\
3 & Demographic & Male sex & Administrative sex M \\
4 & Demographic & Hispanic or Latino & 2135-2 \\
5 & Demographic & Not Hispanic or Latino & 2186-5 \\
6 & Demographic & White & 2106-3 \\
7 & Demographic & Black or African American & 2054-5 \\
8 & Demographic & Asian & 2028-9 \\
9 & Perinatal & Disorders of newborn related to short
gestation and low birth weight & P07 \\
10 & Perinatal & Disorders of newborn related to slow fetal
growth and fetal malnutrition & P05 \\
11 & Perinatal & Newborn affected by maternal factors and
complications of pregnancy, labor, and delivery & P00-P04 \\
12 & Perinatal & Birth trauma & P10-P15 \\
13 & Perinatal & Respiratory distress of newborn & P22 \\
14 & Perinatal & Transitory endocrine and metabolic
disorders specific to newborn & P70-P74 \\
15 & Perinatal & Other disturbances of cerebral status of
newborn & P91 \\
16 & Congenital & Congenital malformations of the nervous
system & Q00-Q07 \\
17 & Congenital & Congenital malformations of the
circulatory system & Q20-Q28 \\
18 & Congenital & Chromosomal abnormalities, not elsewhere
classified & Q90-Q99 \\
19 & Immune & Certain disorders involving the immune
mechanism & D80-D89 \\
20 & Early childhood & Lack of expected normal physiological
development in childhood & R62 \\
21 & Early childhood & Symptoms and signs concerning food
and fluid intake & R63 \\
22 & Early childhood & Acute upper respiratory infections of
multiple and unspecified sites & J06 \\
23 & Early childhood & Asthma & J45 \\
24 & Early childhood & Nonsuppurative otitis media & H65 \\
25 & Early childhood & Suppurative and unspecified otitis
media & H66 \\
26 & Family history & Family history of mental and
behavioral disorders & Z81 \\
27 & Medication & Antibacterials for systemic use & ATC
J01 \\
NA & \emph{Removed (computational limit)} & Feeding problems
of newborn; gastroesophageal reflux disease; family history
of certain disabilities and chronic diseases & P92; K21;
Z82 \\
NA & \emph{Not in model; balance reported} & Female sex;
unknown ethnicity; other, unknown, Native Hawaiian or Other
Pacific Islander, and American Indian or Alaska Native race;
cleft lip and palate; musculoskeletal malformations;
well-child visit & F; UN; 2131-1, UNK, 2076-8, 1002-5;
Q35-Q37; Q65-Q79; Z00.1 \\
\end{longtable}

All diagnosis and medication covariates were binary
indicators assessed at any time before or on the index date.
Codes ending in a 3-character category (eg, P07) include all
subcodes.

\newpage

\hypertarget{etable-3.-outcome-definitions-and-codes}{%
\section*{eTable 3. Outcome Definitions and
Codes}\label{etable-3.-outcome-definitions-and-codes}}

\begin{longtable}[]{@{}
  >{\raggedright\arraybackslash}p{(\columnwidth - 4\tabcolsep) * \real{0.1959}}
  >{\raggedright\arraybackslash}p{(\columnwidth - 4\tabcolsep) * \real{0.4021}}
  >{\raggedright\arraybackslash}p{(\columnwidth - 4\tabcolsep) * \real{0.4021}}@{}}
\toprule\noalign{}
\begin{minipage}[b]{\linewidth}\raggedright
Type
\end{minipage} & \begin{minipage}[b]{\linewidth}\raggedright
Outcome
\end{minipage} & \begin{minipage}[b]{\linewidth}\raggedright
ICD-10-CM codes
\end{minipage} \\
\midrule\noalign{}
\endhead
\bottomrule\noalign{}
\endlastfoot
Primary & Autism spectrum disorder & F84.0 (autistic
disorder), F84.5 (Asperger syndrome), F84.8 (other pervasive
developmental disorders), F84.9 (pervasive developmental
disorder, unspecified) \\
Primary & Attention-deficit/hyperactivity disorder & F90
(all subcodes) \\
Secondary & Speech-language disorder & F80 (specific
developmental disorders of speech and language) \\
Secondary & Intellectual disability & F70-F79 \\
Secondary & Cerebral palsy & G80 \\
Secondary & Epilepsy & G40 \\
Secondary & Sensorineural hearing loss & H90.3, H90.4,
H90.5 \\
Negative control & Fracture of shoulder, upper arm, or
forearm & S42, S52 \\
\end{longtable}

The follow-up window was 1 to 3,285 days after the index
date. Children with the outcome on or before the index date
were excluded from that outcome analysis. F84.2 (Rett
syndrome), F84.3 (other childhood disintegrative disorder),
and F84.4 were not included.

\newpage

\hypertarget{etable-4.-full-results-for-all-conditions-and-outcomes}{%
\section*{eTable 4. Full Results for All Conditions and
Outcomes}\label{etable-4.-full-results-for-all-conditions-and-outcomes}}

\textbf{Autism spectrum disorder}

\begin{longtable}[]{@{}
  >{\raggedright\arraybackslash}p{(\columnwidth - 22\tabcolsep) * \real{0.2348}}
  >{\raggedleft\arraybackslash}p{(\columnwidth - 22\tabcolsep) * \real{0.0606}}
  >{\raggedleft\arraybackslash}p{(\columnwidth - 22\tabcolsep) * \real{0.0606}}
  >{\raggedleft\arraybackslash}p{(\columnwidth - 22\tabcolsep) * \real{0.0606}}
  >{\raggedleft\arraybackslash}p{(\columnwidth - 22\tabcolsep) * \real{0.0606}}
  >{\raggedleft\arraybackslash}p{(\columnwidth - 22\tabcolsep) * \real{0.0606}}
  >{\raggedleft\arraybackslash}p{(\columnwidth - 22\tabcolsep) * \real{0.0606}}
  >{\raggedleft\arraybackslash}p{(\columnwidth - 22\tabcolsep) * \real{0.0606}}
  >{\raggedright\arraybackslash}p{(\columnwidth - 22\tabcolsep) * \real{0.1288}}
  >{\raggedleft\arraybackslash}p{(\columnwidth - 22\tabcolsep) * \real{0.0606}}
  >{\raggedleft\arraybackslash}p{(\columnwidth - 22\tabcolsep) * \real{0.0606}}
  >{\raggedleft\arraybackslash}p{(\columnwidth - 22\tabcolsep) * \real{0.0909}}@{}}
\toprule\noalign{}
\begin{minipage}[b]{\linewidth}\raggedright
Condition
\end{minipage} & \begin{minipage}[b]{\linewidth}\raggedleft
At risk, exposed
\end{minipage} & \begin{minipage}[b]{\linewidth}\raggedleft
Events, exposed
\end{minipage} & \begin{minipage}[b]{\linewidth}\raggedleft
At risk, controls
\end{minipage} & \begin{minipage}[b]{\linewidth}\raggedleft
Events, controls
\end{minipage} & \begin{minipage}[b]{\linewidth}\raggedleft
Risk/1000, exposed
\end{minipage} & \begin{minipage}[b]{\linewidth}\raggedleft
Risk/1000, controls
\end{minipage} & \begin{minipage}[b]{\linewidth}\raggedleft
RD/1000
\end{minipage} & \begin{minipage}[b]{\linewidth}\raggedright
HR (95\% CI)
\end{minipage} & \begin{minipage}[b]{\linewidth}\raggedleft
\emph{P}
\end{minipage} & \begin{minipage}[b]{\linewidth}\raggedleft
FDR \emph{q}
\end{minipage} & \begin{minipage}[b]{\linewidth}\raggedleft
E-value (CI)
\end{minipage} \\
\midrule\noalign{}
\endhead
\bottomrule\noalign{}
\endlastfoot
CMV disease & 1,344 & 108 & 1,342 & 70 & 80.4 & 52.2 & +28.2
& 1.45 (1.07-1.96) & .02 & .03 & 2.25 (1.35) \\
Bacterial meningitis & 1,252 & 81 & 1,253 & 70 & 64.7 & 55.9
& +8.8 & 1.18 (0.86-1.62) & .32 & .44 & 1.64 (1.00) \\
Viral meningitis/encephalitis & 1,685 & 64 & 1,683 & 56 &
38.0 & 33.3 & +4.7 & 1.09 (0.76-1.55) & .65 & .73 & 1.39
(1.00) \\
Sepsis (postneonatal) & 9,650 & 586 & 9,661 & 578 & 60.7 &
59.8 & +0.9 & 0.99 (0.88-1.11) & .85 & .88 & 1.12 (1.00) \\
Osteomyelitis/septic arthritis & 1,261 & 49 & 1,261 & 57 &
38.9 & 45.2 & -6.3 & 0.84 (0.57-1.23) & .36 & .48 & 1.68
(1.00) \\
UTI/pyelonephritis & 57,322 & 1,701 & 57,323 & 1,642 & 29.7
& 28.6 & +1.0 & 0.98 (0.92-1.05) & .56 & .65 & 1.16
(1.00) \\
RSV infection & 149,804 & 5,072 & 149,770 & 5,359 & 33.9 &
35.8 & -1.9 & 0.93 (0.90-0.97) & \textless.001 & .001 & 1.35
(1.21) \\
Influenza & 127,607 & 4,058 & 127,580 & 4,342 & 31.8 & 34.0
& -2.2 & 0.91 (0.87-0.95) & \textless.001 & \textless.001 &
1.43 (1.30) \\
COVID-19 & 94,038 & 4,099 & 94,032 & 3,168 & 43.6 & 33.7 &
+9.9 & 1.22 (1.16-1.28) & \textless.001 & \textless.001 &
1.73 (1.60) \\
Pneumonia & 101,377 & 3,379 & 101,356 & 3,700 & 33.3 & 36.5
& -3.2 & 0.89 (0.84-0.93) & \textless.001 & \textless.001 &
1.51 (1.37) \\
Pertussis & 3,031 & 94 & 3,032 & 117 & 31.0 & 38.6 & -7.6 &
0.77 (0.58-1.00) & .05 & .09 & 1.94 (1.00) \\
HSV infection & 12,697 & 408 & 12,681 & 474 & 32.1 & 37.4 &
-5.2 & 0.82 (0.71-0.93) & .003 & .006 & 1.75 (1.35) \\
Enterovirus infection & 39,879 & 1,668 & 39,874 & 1,735 &
41.8 & 43.5 & -1.7 & 0.93 (0.86-0.99) & .03 & .05 & 1.37
(1.11) \\
Viral gastroenteritis & 112,365 & 4,587 & 112,360 & 4,233 &
40.8 & 37.7 & +3.1 & 1.05 (1.00-1.09) & .03 & .06 & 1.27
(1.06) \\
Roseola & 12,258 & 340 & 12,257 & 330 & 27.7 & 26.9 & +0.8 &
0.97 (0.83-1.13) & .67 & .75 & 1.22 (1.00) \\
Varicella & 4,754 & 146 & 4,751 & 151 & 30.7 & 31.8 & -1.1 &
0.91 (0.72-1.14) & .41 & .53 & 1.43 (1.00) \\
Hand, foot, and mouth disease & 132,498 & 3,665 & 132,435 &
4,624 & 27.7 & 34.9 & -7.3 & 0.77 (0.74-0.80) &
\textless.001 & \textless.001 & 1.93 (1.80) \\
Kawasaki disease & 1,825 & 73 & 1,822 & 74 & 40.0 & 40.6 &
-0.6 & 0.94 (0.68-1.29) & .68 & .75 & 1.34 (1.00) \\
Febrile seizure & 31,532 & 1,353 & 31,526 & 1,242 & 42.9 &
39.4 & +3.5 & 1.05 (0.97-1.14) & .21 & .32 & 1.28 (1.00) \\
Atopic dermatitis & 367,268 & 12,877 & 367,202 & 11,936 &
35.1 & 32.5 & +2.6 & 1.03 (1.01-1.06) & .01 & .03 & 1.21
(1.09) \\
Iron-deficiency anemia & 62,468 & 2,079 & 62,420 & 2,672 &
33.3 & 42.8 & -9.5 & 0.73 (0.69-0.77) & \textless.001 &
\textless.001 & 2.08 (1.91) \\
Diaper dermatitis (reference) & 409,313 & 13,209 & 409,248 &
11,927 & 32.3 & 29.1 & +3.1 & 1.05 (1.03-1.08) &
\textless.001 & NA & 1.29 (1.19) \\
\end{longtable}

\textbf{ADHD}

\begin{longtable}[]{@{}
  >{\raggedright\arraybackslash}p{(\columnwidth - 22\tabcolsep) * \real{0.2348}}
  >{\raggedleft\arraybackslash}p{(\columnwidth - 22\tabcolsep) * \real{0.0606}}
  >{\raggedleft\arraybackslash}p{(\columnwidth - 22\tabcolsep) * \real{0.0606}}
  >{\raggedleft\arraybackslash}p{(\columnwidth - 22\tabcolsep) * \real{0.0606}}
  >{\raggedleft\arraybackslash}p{(\columnwidth - 22\tabcolsep) * \real{0.0606}}
  >{\raggedleft\arraybackslash}p{(\columnwidth - 22\tabcolsep) * \real{0.0606}}
  >{\raggedleft\arraybackslash}p{(\columnwidth - 22\tabcolsep) * \real{0.0606}}
  >{\raggedleft\arraybackslash}p{(\columnwidth - 22\tabcolsep) * \real{0.0606}}
  >{\raggedright\arraybackslash}p{(\columnwidth - 22\tabcolsep) * \real{0.1288}}
  >{\raggedleft\arraybackslash}p{(\columnwidth - 22\tabcolsep) * \real{0.0606}}
  >{\raggedleft\arraybackslash}p{(\columnwidth - 22\tabcolsep) * \real{0.0606}}
  >{\raggedleft\arraybackslash}p{(\columnwidth - 22\tabcolsep) * \real{0.0909}}@{}}
\toprule\noalign{}
\begin{minipage}[b]{\linewidth}\raggedright
Condition
\end{minipage} & \begin{minipage}[b]{\linewidth}\raggedleft
At risk, exposed
\end{minipage} & \begin{minipage}[b]{\linewidth}\raggedleft
Events, exposed
\end{minipage} & \begin{minipage}[b]{\linewidth}\raggedleft
At risk, controls
\end{minipage} & \begin{minipage}[b]{\linewidth}\raggedleft
Events, controls
\end{minipage} & \begin{minipage}[b]{\linewidth}\raggedleft
Risk/1000, exposed
\end{minipage} & \begin{minipage}[b]{\linewidth}\raggedleft
Risk/1000, controls
\end{minipage} & \begin{minipage}[b]{\linewidth}\raggedleft
RD/1000
\end{minipage} & \begin{minipage}[b]{\linewidth}\raggedright
HR (95\% CI)
\end{minipage} & \begin{minipage}[b]{\linewidth}\raggedleft
\emph{P}
\end{minipage} & \begin{minipage}[b]{\linewidth}\raggedleft
FDR \emph{q}
\end{minipage} & \begin{minipage}[b]{\linewidth}\raggedleft
E-value (CI)
\end{minipage} \\
\midrule\noalign{}
\endhead
\bottomrule\noalign{}
\endlastfoot
CMV disease & 1,345 & 77 & 1,345 & 79 & 57.2 & 58.7 & -1.5 &
0.92 (0.67-1.26) & .60 & .68 & 1.40 (1.00) \\
Bacterial meningitis & 1,253 & 61 & 1,253 & 71 & 48.7 & 56.7
& -8.0 & 0.87 (0.62-1.23) & .43 & .53 & 1.56 (1.00) \\
Viral meningitis/encephalitis & 1,685 & 99 & 1,685 & 104 &
58.8 & 61.7 & -3.0 & 0.88 (0.67-1.16) & .38 & .50 & 1.52
(1.00) \\
Sepsis (postneonatal) & 9,661 & 647 & 9,666 & 594 & 67.0 &
61.5 & +5.5 & 1.06 (0.95-1.19) & .27 & .39 & 1.32 (1.00) \\
Osteomyelitis/septic arthritis & 1,262 & 93 & 1,262 & 86 &
73.7 & 68.1 & +5.5 & 1.05 (0.79-1.41) & .73 & .78 & 1.29
(1.00) \\
UTI/pyelonephritis & 57,340 & 2,648 & 57,342 & 2,486 & 46.2
& 43.4 & +2.8 & 1.01 (0.95-1.06) & .80 & .85 & 1.09
(1.00) \\
RSV infection & 149,850 & 7,955 & 149,846 & 7,439 & 53.1 &
49.6 & +3.4 & 1.10 (1.07-1.14) & \textless.001 &
\textless.001 & 1.43 (1.33) \\
Influenza & 127,662 & 6,665 & 127,661 & 6,221 & 52.2 & 48.7
& +3.5 & 1.05 (1.01-1.08) & .009 & .02 & 1.27 (1.12) \\
COVID-19 & 94,098 & 2,413 & 94,099 & 1,966 & 25.6 & 20.9 &
+4.8 & 1.45 (1.36-1.54) & \textless.001 & \textless.001 &
2.25 (2.06) \\
Pneumonia & 101,430 & 6,980 & 101,425 & 6,437 & 68.8 & 63.5
& +5.4 & 1.05 (1.02-1.09) & .003 & .007 & 1.29 (1.15) \\
Pertussis & 3,032 & 204 & 3,033 & 185 & 67.3 & 61.0 & +6.3 &
1.05 (0.86-1.28) & .66 & .74 & 1.27 (1.00) \\
HSV infection & 12,699 & 826 & 12,702 & 772 & 65.0 & 60.8 &
+4.3 & 1.01 (0.92-1.11) & .83 & .86 & 1.12 (1.00) \\
Enterovirus infection & 39,914 & 2,477 & 39,913 & 2,321 &
62.1 & 58.2 & +3.9 & 1.04 (0.98-1.10) & .22 & .32 & 1.23
(1.00) \\
Viral gastroenteritis & 112,427 & 7,827 & 112,430 & 6,282 &
69.6 & 55.9 & +13.7 & 1.21 (1.17-1.25) & \textless.001 &
\textless.001 & 1.70 (1.61) \\
Roseola & 12,264 & 624 & 12,265 & 602 & 50.9 & 49.1 & +1.8 &
0.97 (0.87-1.08) & .58 & .67 & 1.21 (1.00) \\
Varicella & 4,755 & 355 & 4,755 & 305 & 74.7 & 64.1 & +10.5
& 1.08 (0.93-1.26) & .33 & .45 & 1.37 (1.00) \\
Hand, foot, and mouth disease & 132,543 & 7,872 & 132,543 &
7,194 & 59.4 & 54.3 & +5.1 & 1.07 (1.04-1.10) &
\textless.001 & \textless.001 & 1.34 (1.23) \\
Kawasaki disease & 1,825 & 84 & 1,825 & 108 & 46.0 & 59.2 &
-13.2 & 0.73 (0.55-0.97) & .03 & .06 & 2.07 (1.19) \\
Febrile seizure & 31,550 & 2,300 & 31,555 & 1,882 & 72.9 &
59.6 & +13.3 & 1.18 (1.11-1.25) & \textless.001 &
\textless.001 & 1.64 (1.46) \\
Atopic dermatitis & 367,330 & 17,688 & 367,344 & 16,017 &
48.2 & 43.6 & +4.6 & 1.06 (1.04-1.08) & \textless.001 &
\textless.001 & 1.31 (1.23) \\
Iron-deficiency anemia & 62,515 & 3,033 & 62,514 & 3,195 &
48.5 & 51.1 & -2.6 & 0.88 (0.84-0.93) & \textless.001 &
\textless.001 & 1.53 (1.38) \\
Diaper dermatitis (reference) & 409,379 & 25,163 & 409,386 &
17,888 & 61.5 & 43.7 & +17.8 & 1.36 (1.33-1.39) &
\textless.001 & NA & 2.06 (2.00) \\
\end{longtable}

\textbf{Speech-language disorder}

\begin{longtable}[]{@{}
  >{\raggedright\arraybackslash}p{(\columnwidth - 22\tabcolsep) * \real{0.2348}}
  >{\raggedleft\arraybackslash}p{(\columnwidth - 22\tabcolsep) * \real{0.0606}}
  >{\raggedleft\arraybackslash}p{(\columnwidth - 22\tabcolsep) * \real{0.0606}}
  >{\raggedleft\arraybackslash}p{(\columnwidth - 22\tabcolsep) * \real{0.0606}}
  >{\raggedleft\arraybackslash}p{(\columnwidth - 22\tabcolsep) * \real{0.0606}}
  >{\raggedleft\arraybackslash}p{(\columnwidth - 22\tabcolsep) * \real{0.0606}}
  >{\raggedleft\arraybackslash}p{(\columnwidth - 22\tabcolsep) * \real{0.0606}}
  >{\raggedleft\arraybackslash}p{(\columnwidth - 22\tabcolsep) * \real{0.0606}}
  >{\raggedright\arraybackslash}p{(\columnwidth - 22\tabcolsep) * \real{0.1288}}
  >{\raggedleft\arraybackslash}p{(\columnwidth - 22\tabcolsep) * \real{0.0606}}
  >{\raggedleft\arraybackslash}p{(\columnwidth - 22\tabcolsep) * \real{0.0606}}
  >{\raggedleft\arraybackslash}p{(\columnwidth - 22\tabcolsep) * \real{0.0909}}@{}}
\toprule\noalign{}
\begin{minipage}[b]{\linewidth}\raggedright
Condition
\end{minipage} & \begin{minipage}[b]{\linewidth}\raggedleft
At risk, exposed
\end{minipage} & \begin{minipage}[b]{\linewidth}\raggedleft
Events, exposed
\end{minipage} & \begin{minipage}[b]{\linewidth}\raggedleft
At risk, controls
\end{minipage} & \begin{minipage}[b]{\linewidth}\raggedleft
Events, controls
\end{minipage} & \begin{minipage}[b]{\linewidth}\raggedleft
Risk/1000, exposed
\end{minipage} & \begin{minipage}[b]{\linewidth}\raggedleft
Risk/1000, controls
\end{minipage} & \begin{minipage}[b]{\linewidth}\raggedleft
RD/1000
\end{minipage} & \begin{minipage}[b]{\linewidth}\raggedright
HR (95\% CI)
\end{minipage} & \begin{minipage}[b]{\linewidth}\raggedleft
\emph{P}
\end{minipage} & \begin{minipage}[b]{\linewidth}\raggedleft
FDR \emph{q}
\end{minipage} & \begin{minipage}[b]{\linewidth}\raggedleft
E-value (CI)
\end{minipage} \\
\midrule\noalign{}
\endhead
\bottomrule\noalign{}
\endlastfoot
CMV disease & 1,315 & 433 & 1,331 & 343 & 329.3 & 257.7 &
+71.6 & 1.25 (1.08-1.44) & .002 & .005 & 1.80 (1.38) \\
Bacterial meningitis & 1,237 & 358 & 1,245 & 265 & 289.4 &
212.9 & +76.6 & 1.43 (1.22-1.68) & \textless.001 &
\textless.001 & 2.21 (1.74) \\
Viral meningitis/encephalitis & 1,674 & 312 & 1,677 & 314 &
186.4 & 187.2 & -0.9 & 0.94 (0.81-1.10) & .46 & .56 & 1.31
(1.00) \\
Sepsis (postneonatal) & 9,463 & 2,513 & 9,464 & 2,233 &
265.6 & 235.9 & +29.6 & 1.12 (1.06-1.19) & \textless.001 &
\textless.001 & 1.49 (1.31) \\
Osteomyelitis/septic arthritis & 1,230 & 246 & 1,228 & 227 &
200.0 & 184.9 & +15.1 & 1.07 (0.90-1.29) & .43 & .53 & 1.36
(1.00) \\
UTI/pyelonephritis & 56,726 & 8,627 & 56,624 & 8,367 & 152.1
& 147.8 & +4.3 & 0.98 (0.95-1.01) & .21 & .32 & 1.16
(1.00) \\
RSV infection & 147,814 & 26,064 & 147,160 & 24,896 & 176.3
& 169.2 & +7.2 & 1.02 (1.00-1.04) & .02 & .05 & 1.16
(1.06) \\
Influenza & 125,229 & 20,368 & 124,363 & 19,329 & 162.6 &
155.4 & +7.2 & 1.03 (1.01-1.05) & .005 & .01 & 1.20
(1.10) \\
COVID-19 & 92,256 & 19,089 & 91,995 & 14,912 & 206.9 & 162.1
& +44.8 & 1.20 (1.18-1.23) & \textless.001 & \textless.001 &
1.70 (1.63) \\
Pneumonia & 99,196 & 16,670 & 98,812 & 16,316 & 168.1 &
165.1 & +2.9 & 0.99 (0.97-1.01) & .58 & .66 & 1.08 (1.00) \\
Pertussis & 2,997 & 507 & 2,988 & 500 & 169.2 & 167.3 & +1.8
& 0.98 (0.87-1.11) & .81 & .85 & 1.14 (1.00) \\
HSV infection & 12,334 & 2,042 & 12,247 & 1,910 & 165.6 &
156.0 & +9.6 & 1.04 (0.98-1.10) & .24 & .34 & 1.24 (1.00) \\
Enterovirus infection & 38,797 & 7,988 & 38,553 & 7,305 &
205.9 & 189.5 & +16.4 & 1.06 (1.02-1.09) & \textless.001 &
.001 & 1.31 (1.19) \\
Viral gastroenteritis & 109,868 & 21,546 & 109,360 & 18,381
& 196.1 & 168.1 & +28.0 & 1.16 (1.14-1.19) & \textless.001 &
\textless.001 & 1.60 (1.54) \\
Roseola & 12,114 & 1,603 & 12,033 & 1,688 & 132.3 & 140.3 &
-8.0 & 0.89 (0.83-0.96) & .001 & .003 & 1.49 (1.27) \\
Varicella & 4,713 & 686 & 4,646 & 674 & 145.6 & 145.1 & +0.5
& 0.96 (0.86-1.07) & .46 & .56 & 1.25 (1.00) \\
Hand, foot, and mouth disease & 129,145 & 19,723 & 128,271 &
19,938 & 152.7 & 155.4 & -2.7 & 0.96 (0.94-0.98) &
\textless.001 & \textless.001 & 1.25 (1.17) \\
Kawasaki disease & 1,785 & 276 & 1,782 & 289 & 154.6 & 162.2
& -7.6 & 0.90 (0.76-1.06) & .19 & .29 & 1.48 (1.00) \\
Febrile seizure & 30,457 & 5,906 & 30,280 & 4,946 & 193.9 &
163.3 & +30.6 & 1.19 (1.14-1.23) & \textless.001 &
\textless.001 & 1.66 (1.54) \\
Atopic dermatitis & 362,915 & 65,730 & 362,687 & 56,259 &
181.1 & 155.1 & +26.0 & 1.15 (1.13-1.16) & \textless.001 &
\textless.001 & 1.56 (1.52) \\
Iron-deficiency anemia & 60,273 & 11,334 & 59,678 & 10,184 &
188.0 & 170.6 & +17.4 & 1.06 (1.03-1.09) & \textless.001 &
\textless.001 & 1.32 (1.23) \\
Diaper dermatitis (reference) & 404,735 & 70,662 & 403,426 &
57,262 & 174.6 & 141.9 & +32.6 & 1.21 (1.19-1.22) &
\textless.001 & NA & 1.70 (1.67) \\
\end{longtable}

\textbf{Intellectual disability}

\begin{longtable}[]{@{}
  >{\raggedright\arraybackslash}p{(\columnwidth - 22\tabcolsep) * \real{0.2348}}
  >{\raggedleft\arraybackslash}p{(\columnwidth - 22\tabcolsep) * \real{0.0606}}
  >{\raggedleft\arraybackslash}p{(\columnwidth - 22\tabcolsep) * \real{0.0606}}
  >{\raggedleft\arraybackslash}p{(\columnwidth - 22\tabcolsep) * \real{0.0606}}
  >{\raggedleft\arraybackslash}p{(\columnwidth - 22\tabcolsep) * \real{0.0606}}
  >{\raggedleft\arraybackslash}p{(\columnwidth - 22\tabcolsep) * \real{0.0606}}
  >{\raggedleft\arraybackslash}p{(\columnwidth - 22\tabcolsep) * \real{0.0606}}
  >{\raggedleft\arraybackslash}p{(\columnwidth - 22\tabcolsep) * \real{0.0606}}
  >{\raggedright\arraybackslash}p{(\columnwidth - 22\tabcolsep) * \real{0.1288}}
  >{\raggedleft\arraybackslash}p{(\columnwidth - 22\tabcolsep) * \real{0.0606}}
  >{\raggedleft\arraybackslash}p{(\columnwidth - 22\tabcolsep) * \real{0.0606}}
  >{\raggedleft\arraybackslash}p{(\columnwidth - 22\tabcolsep) * \real{0.0909}}@{}}
\toprule\noalign{}
\begin{minipage}[b]{\linewidth}\raggedright
Condition
\end{minipage} & \begin{minipage}[b]{\linewidth}\raggedleft
At risk, exposed
\end{minipage} & \begin{minipage}[b]{\linewidth}\raggedleft
Events, exposed
\end{minipage} & \begin{minipage}[b]{\linewidth}\raggedleft
At risk, controls
\end{minipage} & \begin{minipage}[b]{\linewidth}\raggedleft
Events, controls
\end{minipage} & \begin{minipage}[b]{\linewidth}\raggedleft
Risk/1000, exposed
\end{minipage} & \begin{minipage}[b]{\linewidth}\raggedleft
Risk/1000, controls
\end{minipage} & \begin{minipage}[b]{\linewidth}\raggedleft
RD/1000
\end{minipage} & \begin{minipage}[b]{\linewidth}\raggedright
HR (95\% CI)
\end{minipage} & \begin{minipage}[b]{\linewidth}\raggedleft
\emph{P}
\end{minipage} & \begin{minipage}[b]{\linewidth}\raggedleft
FDR \emph{q}
\end{minipage} & \begin{minipage}[b]{\linewidth}\raggedleft
E-value (CI)
\end{minipage} \\
\midrule\noalign{}
\endhead
\bottomrule\noalign{}
\endlastfoot
CMV disease & 1,345 & 48 & 1,342 & 17 & 35.7 & 12.7 & +23.0
& 2.68 (1.54-4.66) & \textless.001 & .001 & 4.80 (2.45) \\
Bacterial meningitis & 1,248 & 39 & 1,253 & 19 & 31.2 & 15.2
& +16.1 & 2.09 (1.21-3.62) & .008 & .02 & 3.60 (1.71) \\
Viral meningitis/encephalitis & 1,683 & 20 & 1,684 & 11 &
11.9 & 6.5 & +5.4 & 1.71 (0.82-3.57) & .15 & .24 & 2.81
(1.00) \\
Sepsis (postneonatal) & 9,624 & 258 & 9,648 & 206 & 26.8 &
21.4 & +5.5 & 1.23 (1.02-1.48) & .03 & .05 & 1.76 (1.18) \\
Osteomyelitis/septic arthritis & NR & NR & NR & NR & & & &
NR & & & \\
UTI/pyelonephritis & 57,289 & 528 & 57,311 & 443 & 9.2 & 7.7
& +1.5 & 1.13 (1.00-1.28) & .05 & .09 & 1.52 (1.00) \\
RSV infection & 149,813 & 652 & 149,803 & 713 & 4.4 & 4.8 &
-0.4 & 0.93 (0.83-1.03) & .17 & .26 & 1.37 (1.00) \\
Influenza & 127,627 & 402 & 127,608 & 568 & 3.1 & 4.5 & -1.3
& 0.70 (0.61-0.79) & \textless.001 & \textless.001 & 2.23
(1.84) \\
COVID-19 & 94,069 & 212 & 94,046 & 299 & 2.3 & 3.2 & -0.9 &
0.82 (0.68-0.98) & .03 & .06 & 1.74 (1.15) \\
Pneumonia & 101,358 & 904 & 101,367 & 766 & 8.9 & 7.6 & +1.4
& 1.15 (1.04-1.26) & .006 & .01 & 1.56 (1.25) \\
Pertussis & 3,033 & 17 & 3,031 & 20 & 5.6 & 6.6 & -1.0 &
0.80 (0.42-1.53) & .51 & .60 & 1.80 (1.00) \\
HSV infection & 12,699 & 67 & 12,698 & 74 & 5.3 & 5.8 & -0.6
& 0.86 (0.62-1.19) & .36 & .48 & 1.61 (1.00) \\
Enterovirus infection & 39,875 & 317 & 39,868 & 350 & 7.9 &
8.8 & -0.8 & 0.88 (0.75-1.02) & .08 & .14 & 1.55 (1.00) \\
Viral gastroenteritis & 112,385 & 557 & 112,375 & 640 & 5.0
& 5.7 & -0.7 & 0.83 (0.75-0.94) & .002 & .004 & 1.68
(1.34) \\
Roseola & 12,264 & 33 & 12,264 & 41 & 2.7 & 3.3 & -0.7 &
0.76 (0.48-1.20) & .24 & .34 & 1.97 (1.00) \\
Varicella & 4,754 & 21 & 4,752 & 17 & 4.4 & 3.6 & +0.8 &
1.14 (0.60-2.16) & .69 & .75 & 1.54 (1.00) \\
Hand, foot, and mouth disease & 132,552 & 252 & 132,508 &
581 & 1.9 & 4.4 & -2.5 & 0.42 (0.36-0.49) & \textless.001 &
\textless.001 & 4.16 (3.50) \\
Kawasaki disease & 1,824 & 12 & 1,825 & 16 & 6.6 & 8.8 &
-2.2 & 0.73 (0.34-1.54) & .41 & .53 & 2.09 (1.00) \\
Febrile seizure & 31,541 & 245 & 31,537 & 182 & 7.8 & 5.8 &
+2.0 & 1.29 (1.06-1.56) & .009 & .02 & 1.90 (1.32) \\
Atopic dermatitis & 367,331 & 903 & 367,287 & 1,247 & 2.5 &
3.4 & -0.9 & 0.69 (0.64-0.76) & \textless.001 &
\textless.001 & 2.24 (1.98) \\
Iron-deficiency anemia & 62,489 & 274 & 62,465 & 397 & 4.4 &
6.4 & -2.0 & 0.65 (0.56-0.76) & \textless.001 &
\textless.001 & 2.46 (1.98) \\
Diaper dermatitis (reference) & 409,361 & 1,234 & 409,285 &
1,632 & 3.0 & 4.0 & -1.0 & 0.72 (0.67-0.77) & \textless.001
& NA & 2.13 (1.91) \\
\end{longtable}

\textbf{Cerebral palsy}

\begin{longtable}[]{@{}
  >{\raggedright\arraybackslash}p{(\columnwidth - 22\tabcolsep) * \real{0.2348}}
  >{\raggedleft\arraybackslash}p{(\columnwidth - 22\tabcolsep) * \real{0.0606}}
  >{\raggedleft\arraybackslash}p{(\columnwidth - 22\tabcolsep) * \real{0.0606}}
  >{\raggedleft\arraybackslash}p{(\columnwidth - 22\tabcolsep) * \real{0.0606}}
  >{\raggedleft\arraybackslash}p{(\columnwidth - 22\tabcolsep) * \real{0.0606}}
  >{\raggedleft\arraybackslash}p{(\columnwidth - 22\tabcolsep) * \real{0.0606}}
  >{\raggedleft\arraybackslash}p{(\columnwidth - 22\tabcolsep) * \real{0.0606}}
  >{\raggedleft\arraybackslash}p{(\columnwidth - 22\tabcolsep) * \real{0.0606}}
  >{\raggedright\arraybackslash}p{(\columnwidth - 22\tabcolsep) * \real{0.1288}}
  >{\raggedleft\arraybackslash}p{(\columnwidth - 22\tabcolsep) * \real{0.0606}}
  >{\raggedleft\arraybackslash}p{(\columnwidth - 22\tabcolsep) * \real{0.0606}}
  >{\raggedleft\arraybackslash}p{(\columnwidth - 22\tabcolsep) * \real{0.0909}}@{}}
\toprule\noalign{}
\begin{minipage}[b]{\linewidth}\raggedright
Condition
\end{minipage} & \begin{minipage}[b]{\linewidth}\raggedleft
At risk, exposed
\end{minipage} & \begin{minipage}[b]{\linewidth}\raggedleft
Events, exposed
\end{minipage} & \begin{minipage}[b]{\linewidth}\raggedleft
At risk, controls
\end{minipage} & \begin{minipage}[b]{\linewidth}\raggedleft
Events, controls
\end{minipage} & \begin{minipage}[b]{\linewidth}\raggedleft
Risk/1000, exposed
\end{minipage} & \begin{minipage}[b]{\linewidth}\raggedleft
Risk/1000, controls
\end{minipage} & \begin{minipage}[b]{\linewidth}\raggedleft
RD/1000
\end{minipage} & \begin{minipage}[b]{\linewidth}\raggedright
HR (95\% CI)
\end{minipage} & \begin{minipage}[b]{\linewidth}\raggedleft
\emph{P}
\end{minipage} & \begin{minipage}[b]{\linewidth}\raggedleft
FDR \emph{q}
\end{minipage} & \begin{minipage}[b]{\linewidth}\raggedleft
E-value (CI)
\end{minipage} \\
\midrule\noalign{}
\endhead
\bottomrule\noalign{}
\endlastfoot
CMV disease & 1,332 & 115 & 1,334 & 37 & 86.3 & 27.7 & +58.6
& 3.02 (2.09-4.38) & \textless.001 & \textless.001 & 5.49
(3.59) \\
Bacterial meningitis & 1,238 & 174 & 1,246 & 45 & 140.5 &
36.1 & +104.4 & 4.04 (2.91-5.61) & \textless.001 &
\textless.001 & 7.55 (5.27) \\
Viral meningitis/encephalitis & 1,680 & 40 & 1,684 & 15 &
23.8 & 8.9 & +14.9 & 2.58 (1.42-4.66) & .002 & .004 & 4.59
(2.20) \\
Sepsis (postneonatal) & 9,528 & 570 & 9,619 & 318 & 59.8 &
33.1 & +26.8 & 1.80 (1.57-2.06) & \textless.001 &
\textless.001 & 3.00 (2.51) \\
Osteomyelitis/septic arthritis & 1,257 & 29 & 1,260 & 21 &
23.1 & 16.7 & +6.4 & 1.36 (0.78-2.38) & .28 & .40 & 2.06
(1.00) \\
UTI/pyelonephritis & 57,165 & 730 & 57,255 & 673 & 12.8 &
11.8 & +1.0 & 1.04 (0.93-1.15) & .51 & .60 & 1.23 (1.00) \\
RSV infection & 149,673 & 779 & 149,639 & 1,039 & 5.2 & 6.9
& -1.7 & 0.73 (0.67-0.80) & \textless.001 & \textless.001 &
2.08 (1.80) \\
Influenza & 127,556 & 395 & 127,454 & 739 & 3.1 & 5.8 & -2.7
& 0.52 (0.46-0.59) & \textless.001 & \textless.001 & 3.26
(2.80) \\
COVID-19 & 93,987 & 410 & 93,931 & 603 & 4.4 & 6.4 & -2.1 &
0.64 (0.56-0.72) & \textless.001 & \textless.001 & 2.51
(2.11) \\
Pneumonia & 101,120 & 1,225 & 101,240 & 970 & 12.1 & 9.6 &
+2.5 & 1.23 (1.13-1.34) & \textless.001 & \textless.001 &
1.77 (1.52) \\
Pertussis & 3,029 & 17 & 3,026 & 32 & 5.6 & 10.6 & -5.0 &
0.51 (0.28-0.92) & .02 & .05 & 3.34 (1.41) \\
HSV infection & 12,681 & 94 & 12,662 & 80 & 7.4 & 6.3 & +1.1
& 1.13 (0.84-1.52) & .42 & .53 & 1.51 (1.00) \\
Enterovirus infection & 39,751 & 534 & 39,761 & 472 & 13.4 &
11.9 & +1.6 & 1.09 (0.96-1.24) & .16 & .25 & 1.41 (1.00) \\
Viral gastroenteritis & 112,279 & 544 & 112,205 & 832 & 4.8
& 7.4 & -2.6 & 0.63 (0.57-0.70) & \textless.001 &
\textless.001 & 2.54 (2.19) \\
Roseola & 12,256 & 22 & 12,248 & 35 & 1.8 & 2.9 & -1.1 &
0.60 (0.35-1.02) & .06 & .10 & 2.74 (1.00) \\
Varicella & NR & NR & NR & NR & & & & NR & & & \\
Hand, foot, and mouth disease & 132,503 & 212 & 132,324 &
715 & 1.6 & 5.4 & -3.8 & 0.29 (0.25-0.34) & \textless.001 &
\textless.001 & 6.40 (5.40) \\
Kawasaki disease & NR & NR & NR & NR & & & & NR & & & \\
Febrile seizure & 31,492 & 191 & 31,474 & 224 & 6.1 & 7.1 &
-1.1 & 0.83 (0.68-1.00) & .05 & .09 & 1.72 (1.00) \\
Atopic dermatitis & 367,232 & 932 & 367,017 & 1,758 & 2.5 &
4.8 & -2.3 & 0.51 (0.47-0.55) & \textless.001 &
\textless.001 & 3.35 (3.04) \\
Iron-deficiency anemia & 62,437 & 347 & 62,315 & 658 & 5.6 &
10.6 & -5.0 & 0.50 (0.44-0.57) & \textless.001 &
\textless.001 & 3.42 (2.92) \\
Diaper dermatitis (reference) & 409,242 & 1,203 & 408,950 &
2,304 & 2.9 & 5.6 & -2.7 & 0.50 (0.46-0.53) & \textless.001
& NA & 3.44 (3.16) \\
\end{longtable}

\textbf{Epilepsy}

\begin{longtable}[]{@{}
  >{\raggedright\arraybackslash}p{(\columnwidth - 22\tabcolsep) * \real{0.2313}}
  >{\raggedleft\arraybackslash}p{(\columnwidth - 22\tabcolsep) * \real{0.0597}}
  >{\raggedleft\arraybackslash}p{(\columnwidth - 22\tabcolsep) * \real{0.0597}}
  >{\raggedleft\arraybackslash}p{(\columnwidth - 22\tabcolsep) * \real{0.0597}}
  >{\raggedleft\arraybackslash}p{(\columnwidth - 22\tabcolsep) * \real{0.0597}}
  >{\raggedleft\arraybackslash}p{(\columnwidth - 22\tabcolsep) * \real{0.0597}}
  >{\raggedleft\arraybackslash}p{(\columnwidth - 22\tabcolsep) * \real{0.0597}}
  >{\raggedleft\arraybackslash}p{(\columnwidth - 22\tabcolsep) * \real{0.0597}}
  >{\raggedright\arraybackslash}p{(\columnwidth - 22\tabcolsep) * \real{0.1269}}
  >{\raggedleft\arraybackslash}p{(\columnwidth - 22\tabcolsep) * \real{0.0597}}
  >{\raggedleft\arraybackslash}p{(\columnwidth - 22\tabcolsep) * \real{0.0597}}
  >{\raggedleft\arraybackslash}p{(\columnwidth - 22\tabcolsep) * \real{0.1045}}@{}}
\toprule\noalign{}
\begin{minipage}[b]{\linewidth}\raggedright
Condition
\end{minipage} & \begin{minipage}[b]{\linewidth}\raggedleft
At risk, exposed
\end{minipage} & \begin{minipage}[b]{\linewidth}\raggedleft
Events, exposed
\end{minipage} & \begin{minipage}[b]{\linewidth}\raggedleft
At risk, controls
\end{minipage} & \begin{minipage}[b]{\linewidth}\raggedleft
Events, controls
\end{minipage} & \begin{minipage}[b]{\linewidth}\raggedleft
Risk/1000, exposed
\end{minipage} & \begin{minipage}[b]{\linewidth}\raggedleft
Risk/1000, controls
\end{minipage} & \begin{minipage}[b]{\linewidth}\raggedleft
RD/1000
\end{minipage} & \begin{minipage}[b]{\linewidth}\raggedright
HR (95\% CI)
\end{minipage} & \begin{minipage}[b]{\linewidth}\raggedleft
\emph{P}
\end{minipage} & \begin{minipage}[b]{\linewidth}\raggedleft
FDR \emph{q}
\end{minipage} & \begin{minipage}[b]{\linewidth}\raggedleft
E-value (CI)
\end{minipage} \\
\midrule\noalign{}
\endhead
\bottomrule\noalign{}
\endlastfoot
CMV disease & 1,308 & 112 & 1,309 & 47 & 85.6 & 35.9 & +49.7
& 2.34 (1.67-3.30) & \textless.001 & \textless.001 & 4.12
(2.72) \\
Bacterial meningitis & 1,073 & 212 & 1,235 & 59 & 197.6 &
47.8 & +149.8 & 4.45 (3.33-5.93) & \textless.001 &
\textless.001 & 8.36 (6.12) \\
Viral meningitis/encephalitis & 1,606 & 79 & 1,671 & 27 &
49.2 & 16.2 & +33.0 & 2.97 (1.92-4.59) & \textless.001 &
\textless.001 & 5.38 (3.24) \\
Sepsis (postneonatal) & 8,975 & 663 & 9,466 & 337 & 73.9 &
35.6 & +38.3 & 2.07 (1.82-2.36) & \textless.001 &
\textless.001 & 3.56 (3.04) \\
Osteomyelitis/septic arthritis & 1,233 & 39 & 1,241 & 32 &
31.6 & 25.8 & +5.8 & 1.21 (0.76-1.93) & .43 & .53 & 1.71
(1.00) \\
UTI/pyelonephritis & 56,419 & 1,215 & 56,764 & 981 & 21.5 &
17.3 & +4.3 & 1.19 (1.10-1.30) & \textless.001 &
\textless.001 & 1.67 (1.42) \\
RSV infection & 148,929 & 1,840 & 148,889 & 1,935 & 12.4 &
13.0 & -0.6 & 0.94 (0.88-1.00) & .04 & .08 & 1.34 (1.05) \\
Influenza & 126,876 & 1,329 & 126,742 & 1,492 & 10.5 & 11.8
& -1.3 & 0.87 (0.81-0.93) & \textless.001 & \textless.001 &
1.57 (1.35) \\
COVID-19 & 93,419 & 1,067 & 93,405 & 968 & 11.4 & 10.4 &
+1.1 & 1.04 (0.95-1.13) & .40 & .53 & 1.24 (1.00) \\
Pneumonia & 99,987 & 1,879 & 100,439 & 1,596 & 18.8 & 15.9 &
+2.9 & 1.15 (1.08-1.23) & \textless.001 & \textless.001 &
1.57 (1.37) \\
Pertussis & 3,014 & 56 & 3,013 & 54 & 18.6 & 17.9 & +0.7 &
1.00 (0.69-1.45) & .98 & .99 & 1.07 (1.00) \\
HSV infection & 12,571 & 214 & 12,573 & 200 & 17.0 & 15.9 &
+1.1 & 1.03 (0.85-1.25) & .76 & .81 & 1.21 (1.00) \\
Enterovirus infection & 38,990 & 884 & 39,380 & 696 & 22.7 &
17.7 & +5.0 & 1.24 (1.12-1.37) & \textless.001 &
\textless.001 & 1.79 (1.49) \\
Viral gastroenteritis & 111,536 & 1,554 & 111,512 & 1,546 &
13.9 & 13.9 & +0.1 & 0.97 (0.91-1.04) & .43 & .53 & 1.20
(1.00) \\
Roseola & 12,200 & 111 & 12,203 & 122 & 9.1 & 10.0 & -0.9 &
0.86 (0.67-1.11) & .25 & .36 & 1.60 (1.00) \\
Varicella & 4,731 & 38 & 4,731 & 71 & 8.0 & 15.0 & -7.0 &
0.50 (0.34-0.75) & \textless.001 & .002 & 3.38 (2.01) \\
Hand, foot, and mouth disease & 132,081 & 1,008 & 131,570 &
1,555 & 7.6 & 11.8 & -4.2 & 0.63 (0.58-0.68) & \textless.001
& \textless.001 & 2.56 (2.30) \\
Kawasaki disease & 1,815 & 20 & 1,806 & 27 & 11.0 & 15.0 &
-3.9 & 0.70 (0.39-1.24) & .22 & .32 & 2.23 (1.00) \\
Febrile seizure & 30,013 & 2,925 & 31,280 & 377 & 97.5 &
12.1 & +85.4 & 8.33 (7.48-9.28) & \textless.001 &
\textless.001 & 16.15 (14.45) \\
Atopic dermatitis & 366,365 & 3,154 & 365,848 & 4,158 & 8.6
& 11.4 & -2.8 & 0.72 (0.69-0.76) & \textless.001 &
\textless.001 & 2.10 (1.96) \\
Iron-deficiency anemia & 61,937 & 774 & 61,785 & 917 & 12.5
& 14.8 & -2.3 & 0.80 (0.72-0.88) & \textless.001 &
\textless.001 & 1.82 (1.54) \\
Diaper dermatitis (reference) & 408,027 & 4,179 & 407,460 &
4,612 & 10.2 & 11.3 & -1.1 & 0.86 (0.83-0.90) &
\textless.001 & NA & 1.59 (1.46) \\
\end{longtable}

\textbf{Sensorineural hearing loss}

\begin{longtable}[]{@{}
  >{\raggedright\arraybackslash}p{(\columnwidth - 22\tabcolsep) * \real{0.2348}}
  >{\raggedleft\arraybackslash}p{(\columnwidth - 22\tabcolsep) * \real{0.0606}}
  >{\raggedleft\arraybackslash}p{(\columnwidth - 22\tabcolsep) * \real{0.0606}}
  >{\raggedleft\arraybackslash}p{(\columnwidth - 22\tabcolsep) * \real{0.0606}}
  >{\raggedleft\arraybackslash}p{(\columnwidth - 22\tabcolsep) * \real{0.0606}}
  >{\raggedleft\arraybackslash}p{(\columnwidth - 22\tabcolsep) * \real{0.0606}}
  >{\raggedleft\arraybackslash}p{(\columnwidth - 22\tabcolsep) * \real{0.0606}}
  >{\raggedleft\arraybackslash}p{(\columnwidth - 22\tabcolsep) * \real{0.0606}}
  >{\raggedright\arraybackslash}p{(\columnwidth - 22\tabcolsep) * \real{0.1288}}
  >{\raggedleft\arraybackslash}p{(\columnwidth - 22\tabcolsep) * \real{0.0606}}
  >{\raggedleft\arraybackslash}p{(\columnwidth - 22\tabcolsep) * \real{0.0606}}
  >{\raggedleft\arraybackslash}p{(\columnwidth - 22\tabcolsep) * \real{0.0909}}@{}}
\toprule\noalign{}
\begin{minipage}[b]{\linewidth}\raggedright
Condition
\end{minipage} & \begin{minipage}[b]{\linewidth}\raggedleft
At risk, exposed
\end{minipage} & \begin{minipage}[b]{\linewidth}\raggedleft
Events, exposed
\end{minipage} & \begin{minipage}[b]{\linewidth}\raggedleft
At risk, controls
\end{minipage} & \begin{minipage}[b]{\linewidth}\raggedleft
Events, controls
\end{minipage} & \begin{minipage}[b]{\linewidth}\raggedleft
Risk/1000, exposed
\end{minipage} & \begin{minipage}[b]{\linewidth}\raggedleft
Risk/1000, controls
\end{minipage} & \begin{minipage}[b]{\linewidth}\raggedleft
RD/1000
\end{minipage} & \begin{minipage}[b]{\linewidth}\raggedright
HR (95\% CI)
\end{minipage} & \begin{minipage}[b]{\linewidth}\raggedleft
\emph{P}
\end{minipage} & \begin{minipage}[b]{\linewidth}\raggedleft
FDR \emph{q}
\end{minipage} & \begin{minipage}[b]{\linewidth}\raggedleft
E-value (CI)
\end{minipage} \\
\midrule\noalign{}
\endhead
\bottomrule\noalign{}
\endlastfoot
CMV disease & 1,239 & 198 & 1,326 & 43 & 159.8 & 32.4 &
+127.4 & 5.10 (3.67-7.09) & \textless.001 & \textless.001 &
9.67 (6.79) \\
Bacterial meningitis & 1,218 & 120 & 1,242 & 32 & 98.5 &
25.8 & +72.8 & 3.91 (2.64-5.77) & \textless.001 &
\textless.001 & 7.28 (4.73) \\
Viral meningitis/encephalitis & 1,677 & 39 & 1,681 & 22 &
23.3 & 13.1 & +10.2 & 1.70 (1.01-2.87) & .05 & .08 & 2.80
(1.11) \\
Sepsis (postneonatal) & 9,448 & 446 & 9,554 & 291 & 47.2 &
30.5 & +16.7 & 1.53 (1.32-1.78) & \textless.001 &
\textless.001 & 2.44 (1.97) \\
Osteomyelitis/septic arthritis & 1,250 & 33 & 1,256 & 19 &
26.4 & 15.1 & +11.3 & 1.72 (0.97-3.02) & .06 & .10 & 2.82
(1.00) \\
UTI/pyelonephritis & 56,979 & 744 & 57,097 & 683 & 13.1 &
12.0 & +1.1 & 1.04 (0.94-1.16) & .43 & .53 & 1.25 (1.00) \\
RSV infection & 149,302 & 1,232 & 149,314 & 1,502 & 8.3 &
10.1 & -1.8 & 0.81 (0.75-0.87) & \textless.001 &
\textless.001 & 1.79 (1.57) \\
Influenza & 127,214 & 925 & 127,193 & 1,136 & 7.3 & 8.9 &
-1.7 & 0.79 (0.73-0.86) & \textless.001 & \textless.001 &
1.83 (1.58) \\
COVID-19 & 93,655 & 785 & 93,715 & 847 & 8.4 & 9.0 & -0.7 &
0.88 (0.80-0.97) & .01 & .03 & 1.52 (1.19) \\
Pneumonia & 100,800 & 1,257 & 100,982 & 1,156 & 12.5 & 11.4
& +1.0 & 1.06 (0.98-1.15) & .15 & .24 & 1.31 (1.00) \\
Pertussis & 3,016 & 30 & 3,022 & 32 & 9.9 & 10.6 & -0.6 &
0.90 (0.55-1.49) & .69 & .75 & 1.46 (1.00) \\
HSV infection & 12,644 & 111 & 12,647 & 121 & 8.8 & 9.6 &
-0.8 & 0.88 (0.68-1.14) & .33 & .45 & 1.53 (1.00) \\
Enterovirus infection & 39,558 & 540 & 39,640 & 523 & 13.7 &
13.2 & +0.5 & 1.00 (0.89-1.13) & 1.00 & 1.00 & 1.00
(1.00) \\
Viral gastroenteritis & 111,858 & 1,169 & 111,998 & 1,140 &
10.5 & 10.2 & +0.3 & 0.99 (0.92-1.08) & .87 & .88 & 1.09
(1.00) \\
Roseola & 12,242 & 89 & 12,241 & 82 & 7.3 & 6.7 & +0.6 &
1.03 (0.76-1.39) & .86 & .88 & 1.20 (1.00) \\
Varicella & 4,743 & 36 & 4,746 & 29 & 7.6 & 6.1 & +1.5 &
1.18 (0.72-1.92) & .51 & .60 & 1.64 (1.00) \\
Hand, foot, and mouth disease & 132,033 & 864 & 132,075 &
1,063 & 6.5 & 8.0 & -1.5 & 0.79 (0.72-0.87) & \textless.001
& \textless.001 & 1.84 (1.58) \\
Kawasaki disease & 1,823 & 14 & 1,823 & 14 & 7.7 & 7.7 &
+0.0 & 0.95 (0.45-2.00) & .89 & .91 & 1.28 (1.00) \\
Febrile seizure & 31,386 & 269 & 31,404 & 298 & 8.6 & 9.5 &
-0.9 & 0.87 (0.74-1.03) & .11 & .17 & 1.55 (1.00) \\
Atopic dermatitis & 366,334 & 3,099 & 366,462 & 3,206 & 8.5
& 8.7 & -0.3 & 0.93 (0.88-0.97) & .003 & .006 & 1.37
(1.19) \\
Iron-deficiency anemia & 62,176 & 595 & 62,158 & 661 & 9.6 &
10.6 & -1.1 & 0.85 (0.76-0.95) & .005 & .01 & 1.62 (1.27) \\
Diaper dermatitis (reference) & 408,275 & 3,290 & 408,313 &
3,581 & 8.1 & 8.8 & -0.7 & 0.88 (0.84-0.92) & \textless.001
& NA & 1.54 (1.40) \\
\end{longtable}

\textbf{Arm fracture (negative-control outcome)}

\begin{longtable}[]{@{}
  >{\raggedright\arraybackslash}p{(\columnwidth - 22\tabcolsep) * \real{0.2348}}
  >{\raggedleft\arraybackslash}p{(\columnwidth - 22\tabcolsep) * \real{0.0606}}
  >{\raggedleft\arraybackslash}p{(\columnwidth - 22\tabcolsep) * \real{0.0606}}
  >{\raggedleft\arraybackslash}p{(\columnwidth - 22\tabcolsep) * \real{0.0606}}
  >{\raggedleft\arraybackslash}p{(\columnwidth - 22\tabcolsep) * \real{0.0606}}
  >{\raggedleft\arraybackslash}p{(\columnwidth - 22\tabcolsep) * \real{0.0606}}
  >{\raggedleft\arraybackslash}p{(\columnwidth - 22\tabcolsep) * \real{0.0606}}
  >{\raggedleft\arraybackslash}p{(\columnwidth - 22\tabcolsep) * \real{0.0606}}
  >{\raggedright\arraybackslash}p{(\columnwidth - 22\tabcolsep) * \real{0.1288}}
  >{\raggedleft\arraybackslash}p{(\columnwidth - 22\tabcolsep) * \real{0.0606}}
  >{\raggedleft\arraybackslash}p{(\columnwidth - 22\tabcolsep) * \real{0.0606}}
  >{\raggedleft\arraybackslash}p{(\columnwidth - 22\tabcolsep) * \real{0.0909}}@{}}
\toprule\noalign{}
\begin{minipage}[b]{\linewidth}\raggedright
Condition
\end{minipage} & \begin{minipage}[b]{\linewidth}\raggedleft
At risk, exposed
\end{minipage} & \begin{minipage}[b]{\linewidth}\raggedleft
Events, exposed
\end{minipage} & \begin{minipage}[b]{\linewidth}\raggedleft
At risk, controls
\end{minipage} & \begin{minipage}[b]{\linewidth}\raggedleft
Events, controls
\end{minipage} & \begin{minipage}[b]{\linewidth}\raggedleft
Risk/1000, exposed
\end{minipage} & \begin{minipage}[b]{\linewidth}\raggedleft
Risk/1000, controls
\end{minipage} & \begin{minipage}[b]{\linewidth}\raggedleft
RD/1000
\end{minipage} & \begin{minipage}[b]{\linewidth}\raggedright
HR (95\% CI)
\end{minipage} & \begin{minipage}[b]{\linewidth}\raggedleft
\emph{P}
\end{minipage} & \begin{minipage}[b]{\linewidth}\raggedleft
FDR \emph{q}
\end{minipage} & \begin{minipage}[b]{\linewidth}\raggedleft
E-value (CI)
\end{minipage} \\
\midrule\noalign{}
\endhead
\bottomrule\noalign{}
\endlastfoot
CMV disease & 1,334 & 37 & 1,336 & 40 & 27.7 & 29.9 & -2.2 &
0.86 (0.55-1.35) & .51 & NA & 1.59 (1.00) \\
Bacterial meningitis & 1,241 & 40 & 1,248 & 34 & 32.2 & 27.2
& +5.0 & 1.20 (0.76-1.89) & .45 & NA & 1.68 (1.00) \\
Viral meningitis/encephalitis & 1,678 & 58 & 1,683 & 63 &
34.6 & 37.4 & -2.9 & 0.85 (0.60-1.22) & .39 & NA & 1.61
(1.00) \\
Sepsis (postneonatal) & 9,582 & 325 & 9,619 & 273 & 33.9 &
28.4 & +5.5 & 1.17 (0.99-1.37) & .06 & NA & 1.61 (1.00) \\
Osteomyelitis/septic arthritis & 1,223 & 72 & 1,260 & 23 &
58.9 & 18.3 & +40.6 & 3.27 (2.04-5.22) & \textless.001 & NA
& 5.98 (3.50) \\
UTI/pyelonephritis & 57,147 & 2,140 & 57,210 & 1,806 & 37.4
& 31.6 & +5.9 & 1.13 (1.06-1.20) & \textless.001 & NA & 1.51
(1.32) \\
RSV infection & 149,485 & 4,784 & 149,515 & 4,320 & 32.0 &
28.9 & +3.1 & 1.11 (1.07-1.16) & \textless.001 & NA & 1.47
(1.34) \\
Influenza & 127,265 & 4,176 & 127,332 & 3,538 & 32.8 & 27.8
& +5.0 & 1.16 (1.11-1.21) & \textless.001 & NA & 1.59
(1.46) \\
COVID-19 & 93,813 & 1,779 & 93,862 & 1,474 & 19.0 & 15.7 &
+3.3 & 1.20 (1.12-1.28) & \textless.001 & NA & 1.69
(1.48) \\
Pneumonia & 101,074 & 3,911 & 101,137 & 3,460 & 38.7 & 34.2
& +4.5 & 1.10 (1.05-1.15) & \textless.001 & NA & 1.43
(1.29) \\
Pertussis & 3,028 & 123 & 3,026 & 96 & 40.6 & 31.7 & +8.9 &
1.23 (0.94-1.61) & .13 & NA & 1.77 (1.00) \\
HSV infection & 12,625 & 517 & 12,660 & 451 & 41.0 & 35.6 &
+5.3 & 1.10 (0.97-1.25) & .15 & NA & 1.42 (1.00) \\
Enterovirus infection & 39,743 & 1,401 & 39,803 & 1,206 &
35.3 & 30.3 & +5.0 & 1.13 (1.05-1.22) & .002 & NA & 1.52
(1.27) \\
Viral gastroenteritis & 112,019 & 4,348 & 112,139 & 3,660 &
38.8 & 32.6 & +6.2 & 1.15 (1.10-1.20) & \textless.001 & NA &
1.56 (1.43) \\
Roseola & 12,244 & 421 & 12,241 & 390 & 34.4 & 31.9 & +2.5 &
1.02 (0.89-1.17) & .81 & NA & 1.15 (1.00) \\
Varicella & 4,742 & 215 & 4,748 & 184 & 45.3 & 38.8 & +6.6 &
1.09 (0.90-1.33) & .37 & NA & 1.41 (1.00) \\
Hand, foot, and mouth disease & 132,097 & 4,622 & 132,179 &
4,121 & 35.0 & 31.2 & +3.8 & 1.10 (1.05-1.15) &
\textless.001 & NA & 1.43 (1.29) \\
Kawasaki disease & 1,819 & 59 & 1,820 & 45 & 32.4 & 24.7 &
+7.7 & 1.26 (0.85-1.85) & .25 & NA & 1.83 (1.00) \\
Febrile seizure & 31,438 & 1,181 & 31,459 & 1,034 & 37.6 &
32.9 & +4.7 & 1.10 (1.01-1.20) & .03 & NA & 1.43 (1.12) \\
Atopic dermatitis & 366,586 & 10,297 & 366,675 & 9,484 &
28.1 & 25.9 & +2.2 & 1.04 (1.01-1.07) & .005 & NA & 1.25
(1.12) \\
Iron-deficiency anemia & 62,281 & 1,905 & 62,314 & 1,719 &
30.6 & 27.6 & +3.0 & 1.05 (0.98-1.12) & .18 & NA & 1.27
(1.00) \\
Diaper dermatitis (reference) & 408,509 & 13,281 & 408,557 &
11,831 & 32.5 & 29.0 & +3.6 & 1.07 (1.04-1.10) &
\textless.001 & NA & 1.34 (1.25) \\
\end{longtable}

Abbreviations: NA, not applicable or not computed; ADHD,
attention-deficit/hyperactivity disorder; CMV,
cytomegalovirus; FDR, false discovery rate; HR, hazard
ratio; HSV, herpes simplex virus; NR, not reported
(suppressed count); RD, risk difference; RSV, respiratory
syncytial virus; UTI, urinary tract infection. Risks are
crude cumulative proportions per 1,000 children over 1 day
to 9 years. \emph{P} values were derived from the HR and
95\% CI. FDR \emph{q} values are Benjamini-Hochberg adjusted
across 144 tests; not computed (NA) for the reference
exposure (diaper dermatitis) or the negative-control
outcome. E-values are for the point estimate and (in
parentheses) the CI limit closest to 1. Event counts of 10
or fewer are shown as ≤10.

\newpage

\hypertarget{etable-5.-hazard-ratios-calibrated-against-diaper-dermatitis}{%
\section*{eTable 5. Hazard Ratios Calibrated Against Diaper
Dermatitis}\label{etable-5.-hazard-ratios-calibrated-against-diaper-dermatitis}}

\begin{longtable}[]{@{}
  >{\raggedright\arraybackslash}p{(\columnwidth - 14\tabcolsep) * \real{0.2092}}
  >{\raggedleft\arraybackslash}p{(\columnwidth - 14\tabcolsep) * \real{0.1111}}
  >{\raggedleft\arraybackslash}p{(\columnwidth - 14\tabcolsep) * \real{0.1111}}
  >{\raggedleft\arraybackslash}p{(\columnwidth - 14\tabcolsep) * \real{0.1111}}
  >{\raggedleft\arraybackslash}p{(\columnwidth - 14\tabcolsep) * \real{0.1111}}
  >{\raggedleft\arraybackslash}p{(\columnwidth - 14\tabcolsep) * \real{0.1176}}
  >{\raggedleft\arraybackslash}p{(\columnwidth - 14\tabcolsep) * \real{0.1176}}
  >{\raggedleft\arraybackslash}p{(\columnwidth - 14\tabcolsep) * \real{0.1111}}@{}}
\toprule\noalign{}
\begin{minipage}[b]{\linewidth}\raggedright
Condition
\end{minipage} & \begin{minipage}[b]{\linewidth}\raggedleft
Autism
\end{minipage} & \begin{minipage}[b]{\linewidth}\raggedleft
ADHD
\end{minipage} & \begin{minipage}[b]{\linewidth}\raggedleft
Speech-language
\end{minipage} & \begin{minipage}[b]{\linewidth}\raggedleft
Intellectual disability
\end{minipage} & \begin{minipage}[b]{\linewidth}\raggedleft
Cerebral palsy
\end{minipage} & \begin{minipage}[b]{\linewidth}\raggedleft
Epilepsy
\end{minipage} & \begin{minipage}[b]{\linewidth}\raggedleft
Hearing loss
\end{minipage} \\
\midrule\noalign{}
\endhead
\bottomrule\noalign{}
\endlastfoot
CMV disease & 1.38 (1.02-1.86) & 0.68 (0.49-0.93) & 1.04
(0.90-1.19) & 3.73 (2.13-6.52) & 6.08 (4.17-8.86) & 2.72
(1.93-3.83) & 5.81 (4.17-8.11) \\
Bacterial meningitis & 1.12 (0.81-1.54) & 0.64 (0.46-0.90) &
1.19 (1.01-1.39) & 2.91 (1.67-5.06) & 8.13 (5.82-11.37) &
5.15 (3.85-6.90) & 4.45 (3.01-6.60) \\
Viral meningitis/encephalitis & 1.03 (0.72-1.48) & 0.65
(0.49-0.86) & 0.78 (0.67-0.92) & 2.38 (1.14-4.99) & 5.18
(2.85-9.42) & 3.44 (2.22-5.33) & 1.94 (1.15-3.28) \\
Sepsis (postneonatal) & 0.94 (0.84-1.06) & 0.78 (0.70-0.88)
& 0.93 (0.88-0.99) & 1.71 (1.40-2.09) & 3.62 (3.10-4.22) &
2.40 (2.09-2.75) & 1.75 (1.50-2.04) \\
Osteomyelitis/septic arthritis & 0.79 (0.54-1.16) & 0.78
(0.58-1.04) & 0.89 (0.74-1.07) & NR & 2.73 (1.55-4.82) &
1.40 (0.87-2.24) & 1.96 (1.11-3.45) \\
UTI/pyelonephritis & 0.93 (0.87-1.00) & 0.74 (0.70-0.79) &
0.81 (0.79-0.84) & 1.58 (1.36-1.83) & 2.08 (1.84-2.36) &
1.38 (1.26-1.52) & 1.19 (1.06-1.33) \\
RSV infection & 0.89 (0.85-0.93) & 0.81 (0.78-0.84) & 0.85
(0.83-0.86) & 1.29 (1.14-1.47) & 1.47 (1.31-1.65) & 1.08
(1.00-1.17) & 0.92 (0.84-1.00) \\
Influenza & 0.86 (0.82-0.91) & 0.77 (0.74-0.80) & 0.85
(0.83-0.87) & 0.97 (0.84-1.12) & 1.05 (0.91-1.20) & 1.00
(0.92-1.09) & 0.90 (0.82-1.00) \\
COVID-19 & 1.16 (1.10-1.22) & 1.07 (1.00-1.14) & 1.00
(0.97-1.02) & 1.14 (0.93-1.39) & 1.28 (1.11-1.48) & 1.20
(1.09-1.33) & 1.01 (0.90-1.12) \\
Pneumonia & 0.84 (0.80-0.89) & 0.77 (0.75-0.81) & 0.82
(0.81-0.84) & 1.60 (1.41-1.80) & 2.48 (2.22-2.76) & 1.33
(1.23-1.44) & 1.21 (1.10-1.33) \\
Pertussis & 0.73 (0.55-0.95) & 0.77 (0.63-0.94) & 0.82
(0.72-0.93) & 1.12 (0.58-2.14) & 1.02 (0.57-1.85) & 1.15
(0.79-1.68) & 1.03 (0.62-1.70) \\
HSV infection & 0.77 (0.68-0.89) & 0.74 (0.67-0.82) & 0.86
(0.81-0.92) & 1.19 (0.85-1.67) & 2.27 (1.67-3.08) & 1.19
(0.98-1.46) & 1.00 (0.77-1.30) \\
Enterovirus infection & 0.88 (0.82-0.94) & 0.76 (0.72-0.81)
& 0.88 (0.85-0.91) & 1.22 (1.03-1.44) & 2.20 (1.91-2.54) &
1.44 (1.29-1.60) & 1.14 (1.00-1.30) \\
Viral gastroenteritis & 0.99 (0.95-1.04) & 0.89 (0.85-0.92)
& 0.96 (0.94-0.99) & 1.16 (1.02-1.33) & 1.27 (1.12-1.45) &
1.13 (1.04-1.22) & 1.13 (1.03-1.24) \\
Roseola & 0.92 (0.79-1.07) & 0.71 (0.64-0.80) & 0.74
(0.69-0.79) & 1.06 (0.66-1.68) & 1.20 (0.70-2.05) & 1.00
(0.77-1.29) & 1.17 (0.87-1.59) \\
Varicella & 0.86 (0.69-1.09) & 0.79 (0.68-0.93) & 0.80
(0.72-0.89) & 1.59 (0.83-3.02) & NR & 0.58 (0.39-0.87) &
1.34 (0.82-2.20) \\
Hand, foot, and mouth disease & 0.73 (0.69-0.77) & 0.79
(0.76-0.82) & 0.80 (0.78-0.82) & 0.59 (0.50-0.69) & 0.58
(0.49-0.69) & 0.73 (0.67-0.80) & 0.90 (0.81-1.00) \\
Kawasaki disease & 0.89 (0.64-1.23) & 0.54 (0.41-0.72) &
0.74 (0.63-0.88) & 1.01 (0.48-2.15) & NR & 0.81 (0.45-1.44)
& 1.08 (0.52-2.28) \\
Febrile seizure & 1.00 (0.92-1.08) & 0.87 (0.81-0.92) & 0.98
(0.95-1.02) & 1.80 (1.46-2.21) & 1.66 (1.35-2.04) & 9.65
(8.60-10.83) & 1.00 (0.84-1.18) \\
Atopic dermatitis & 0.98 (0.95-1.01) & 0.78 (0.76-0.80) &
0.95 (0.94-0.97) & 0.97 (0.86-1.08) & 1.02 (0.92-1.14) &
0.84 (0.79-0.89) & 1.06 (0.99-1.13) \\
Iron-deficiency anemia & 0.69 (0.65-0.74) & 0.65 (0.61-0.68)
& 0.88 (0.86-0.91) & 0.90 (0.76-1.07) & 1.00 (0.87-1.16) &
0.92 (0.83-1.02) & 0.97 (0.86-1.10) \\
\emph{Reference HR, diaper dermatitis} & \emph{1.05
(1.03-1.08)} & \emph{1.36 (1.33-1.39)} & \emph{1.21
(1.19-1.22)} & \emph{0.72 (0.67-0.77)} & \emph{0.50
(0.46-0.53)} & \emph{0.86 (0.83-0.90)} & \emph{0.88
(0.84-0.92)} \\
\end{longtable}

Values are calibrated HR (95\% CI) = HR for the condition
divided by the HR for diaper dermatitis for the same
outcome; see eMethods 6. Calibrated HRs above 1 indicate
risk beyond that associated with a benign sick visit.

\newpage

\hypertarget{etable-6.-random-effects-pooled-hazard-ratios-by-condition-group}{%
\section*{eTable 6. Random-Effects Pooled Hazard Ratios by
Condition
Group}\label{etable-6.-random-effects-pooled-hazard-ratios-by-condition-group}}

\begin{longtable}[]{@{}
  >{\raggedright\arraybackslash}p{(\columnwidth - 6\tabcolsep) * \real{0.3535}}
  >{\raggedright\arraybackslash}p{(\columnwidth - 6\tabcolsep) * \real{0.3939}}
  >{\raggedright\arraybackslash}p{(\columnwidth - 6\tabcolsep) * \real{0.1717}}
  >{\raggedleft\arraybackslash}p{(\columnwidth - 6\tabcolsep) * \real{0.0808}}@{}}
\toprule\noalign{}
\begin{minipage}[b]{\linewidth}\raggedright
Condition group (k)
\end{minipage} & \begin{minipage}[b]{\linewidth}\raggedright
Outcome
\end{minipage} & \begin{minipage}[b]{\linewidth}\raggedright
Pooled HR (95\% CI)
\end{minipage} & \begin{minipage}[b]{\linewidth}\raggedleft
\emph{I}\textsuperscript{2}, \%
\end{minipage} \\
\midrule\noalign{}
\endhead
\bottomrule\noalign{}
\endlastfoot
Invasive CNS (2) & Autism spectrum disorder & 1.14
(0.90-1.44) & 0 \\
Invasive CNS (2) & ADHD & 0.88 (0.71-1.09) & 0 \\
Invasive CNS (2) & Speech-language disorder & 1.16
(0.77-1.74) & 92 \\
Invasive CNS (2) & Intellectual disability & 1.94
(1.25-3.02) & 0 \\
Invasive CNS (2) & Cerebral palsy & 3.46 (2.28-5.27) & 41 \\
Invasive CNS (2) & Epilepsy & 3.76 (2.54-5.55) & 56 \\
Invasive CNS (2) & Sensorineural hearing loss & 2.63
(1.17-5.92) & 84 \\
Invasive CNS (2) & Arm fracture (negative-control outcome) &
0.98 (0.71-1.35) & 22 \\
Invasive systemic (3) & Autism spectrum disorder & 0.98
(0.92-1.04) & 0 \\
Invasive systemic (3) & ADHD & 1.02 (0.97-1.07) & 0 \\
Invasive systemic (3) & Speech-language disorder & 1.05
(0.94-1.17) & 88 \\
Invasive systemic (3) & Intellectual disability & 1.16
(1.05-1.29) & 0 \\
Invasive systemic (3) & Cerebral palsy & 1.36 (0.87-2.13) &
95 \\
Invasive systemic (3) & Epilepsy & 1.46 (0.94-2.28) & 96 \\
Invasive systemic (3) & Sensorineural hearing loss & 1.34
(0.96-1.86) & 90 \\
Invasive systemic (3) & Arm fracture (negative-control
outcome) & 1.44 (1.05-1.96) & 90 \\
Respiratory (5) & Autism spectrum disorder & 0.95
(0.84-1.08) & 97 \\
Respiratory (5) & ADHD & 1.13 (1.03-1.25) & 96 \\
Respiratory (5) & Speech-language disorder & 1.05
(0.97-1.13) & 98 \\
Respiratory (5) & Intellectual disability & 0.88 (0.71-1.09)
& 90 \\
Respiratory (5) & Cerebral palsy & 0.70 (0.49-1.00) & 98 \\
Respiratory (5) & Epilepsy & 0.99 (0.88-1.12) & 89 \\
Respiratory (5) & Sensorineural hearing loss & 0.88
(0.77-1.00) & 87 \\
Respiratory (5) & Arm fracture (negative-control outcome) &
1.14 (1.10-1.17) & 33 \\
Mucocutaneous/enteric/exanthem (6) & Autism spectrum
disorder & 0.90 (0.78-1.03) & 95 \\
Mucocutaneous/enteric/exanthem (6) & ADHD & 1.07 (0.99-1.14)
& 89 \\
Mucocutaneous/enteric/exanthem (6) & Speech-language
disorder & 1.01 (0.93-1.11) & 98 \\
Mucocutaneous/enteric/exanthem (6) & Intellectual disability
& 0.76 (0.55-1.03) & 92 \\
Mucocutaneous/enteric/exanthem (6) & Cerebral palsy & 0.67
(0.40-1.13) & 98 \\
Mucocutaneous/enteric/exanthem (6) & Epilepsy & 0.85
(0.66-1.10) & 96 \\
Mucocutaneous/enteric/exanthem (6) & Sensorineural hearing
loss & 0.94 (0.83-1.06) & 72 \\
Mucocutaneous/enteric/exanthem (6) & Arm fracture
(negative-control outcome) & 1.12 (1.09-1.15) & 0 \\
Noninfectious/inflammatory (4) & Autism spectrum disorder &
0.93 (0.75-1.14) & 98 \\
Noninfectious/inflammatory (4) & ADHD & 0.99 (0.87-1.13) &
96 \\
Noninfectious/inflammatory (4) & Speech-language disorder &
1.11 (1.05-1.17) & 92 \\
Noninfectious/inflammatory (4) & Intellectual disability &
0.81 (0.59-1.12) & 92 \\
Noninfectious/inflammatory (4) & Cerebral palsy & 0.59
(0.46-0.75) & 91 \\
Noninfectious/inflammatory (4) & Epilepsy & 1.37 (0.42-4.44)
& 100 \\
Noninfectious/inflammatory (4) & Sensorineural hearing loss
& 0.91 (0.87-0.95) & 0 \\
Noninfectious/inflammatory (4) & Arm fracture
(negative-control outcome) & 1.05 (1.02-1.07) & 0 \\
\end{longtable}

DerSimonian-Laird random-effects pooling of log HRs.
Invasive CNS: bacterial meningitis, viral
meningitis/encephalitis. Invasive systemic: postneonatal
sepsis, UTI/pyelonephritis, osteomyelitis/septic arthritis.
Respiratory: RSV, influenza, COVID-19, pneumonia, pertussis.
Mucocutaneous, enteric, and exanthematous: HSV, enterovirus,
viral gastroenteritis, roseola, varicella, hand, foot, and
mouth disease. Noninfectious/inflammatory: Kawasaki disease,
febrile seizure, atopic dermatitis, iron-deficiency anemia.
CMV was not pooled.

\newpage

\hypertarget{etable-7.-sensitivity-analysis-without-the-well-child-requirement-in-exposed-infants}{%
\section*{eTable 7. Sensitivity Analysis Without the
Well-Child Requirement in Exposed
Infants}\label{etable-7.-sensitivity-analysis-without-the-well-child-requirement-in-exposed-infants}}

\begin{longtable}[]{@{}
  >{\raggedright\arraybackslash}p{(\columnwidth - 8\tabcolsep) * \real{0.2870}}
  >{\raggedright\arraybackslash}p{(\columnwidth - 8\tabcolsep) * \real{0.1667}}
  >{\raggedright\arraybackslash}p{(\columnwidth - 8\tabcolsep) * \real{0.2315}}
  >{\raggedright\arraybackslash}p{(\columnwidth - 8\tabcolsep) * \real{0.1574}}
  >{\raggedright\arraybackslash}p{(\columnwidth - 8\tabcolsep) * \real{0.1574}}@{}}
\toprule\noalign{}
\begin{minipage}[b]{\linewidth}\raggedright
Condition
\end{minipage} & \begin{minipage}[b]{\linewidth}\raggedright
Matched pairs, unrestricted / primary
\end{minipage} & \begin{minipage}[b]{\linewidth}\raggedright
Outcome
\end{minipage} & \begin{minipage}[b]{\linewidth}\raggedright
HR (95\% CI), unrestricted
\end{minipage} & \begin{minipage}[b]{\linewidth}\raggedright
HR (95\% CI), primary
\end{minipage} \\
\midrule\noalign{}
\endhead
\bottomrule\noalign{}
\endlastfoot
CMV disease & 5,147 / 1,345 & Autism & 1.45 (1.21-1.73) &
1.45 (1.07-1.96) \\
& & ADHD & 0.87 (0.72-1.05) & 0.92 (0.67-1.26) \\
Bacterial meningitis & 6,724 / 1,253 & Autism & 1.14
(0.95-1.36) & 1.18 (0.86-1.62) \\
& & ADHD & 0.62 (0.52-0.73) & 0.87 (0.62-1.23) \\
& & Speech-language disorder & 0.95 (0.87-1.03) & 1.43
(1.22-1.68) \\
& & Intellectual disability & 1.77 (1.35-2.32) & 2.09
(1.21-3.62) \\
& & Cerebral palsy & 3.50 (2.96-4.14) & 4.04 (2.91-5.61) \\
& & Epilepsy & 5.78 (4.94-6.77) & 4.45 (3.33-5.93) \\
& & Hearing loss & 3.85 (3.18-4.67) & 3.91 (2.64-5.77) \\
& & Arm fracture & 0.83 (0.67-1.02) & 1.20 (0.76-1.89) \\
Viral meningitis/encephalitis & 8,707 / 1,685 & Autism &
0.91 (0.76-1.10) & 1.09 (0.76-1.55) \\
& & ADHD & 0.68 (0.58-0.79) & 0.88 (0.67-1.16) \\
& & Speech-language disorder & 0.64 (0.58-0.70) & 0.94
(0.81-1.10) \\
& & Intellectual disability & 1.77 (1.14-2.75) & 1.71
(0.82-3.57) \\
& & Cerebral palsy & 3.10 (2.24-4.31) & 2.58 (1.42-4.66) \\
& & Epilepsy & 3.72 (2.98-4.63) & 2.97 (1.92-4.59) \\
& & Hearing loss & 1.68 (1.24-2.27) & 1.70 (1.01-2.87) \\
& & Arm fracture & 1.05 (0.89-1.24) & 0.85 (0.60-1.22) \\
Sepsis (postneonatal) & 51,429 / 9,667 & Autism & 0.92
(0.86-0.98) & 0.99 (0.88-1.11) \\
& & ADHD & 0.77 (0.72-0.81) & 1.06 (0.95-1.19) \\
& & Speech-language disorder & 0.71 (0.68-0.73) & 1.12
(1.06-1.19) \\
& & Intellectual disability & 1.45 (1.32-1.60) & 1.23
(1.02-1.48) \\
& & Cerebral palsy & 1.80 (1.67-1.94) & 1.80 (1.57-2.06) \\
& & Epilepsy & 2.28 (2.13-2.44) & 2.07 (1.82-2.36) \\
& & Hearing loss & 1.46 (1.35-1.59) & 1.53 (1.32-1.78) \\
& & Arm fracture & 1.03 (0.95-1.11) & 1.17 (0.99-1.37) \\
Osteomyelitis/septic arthritis & 7,718 / 1,262 & Autism &
0.95 (0.78-1.16) & 0.84 (0.57-1.23) \\
& & ADHD & 0.86 (0.74-1.00) & 1.05 (0.79-1.41) \\
& & Speech-language disorder & 0.67 (0.60-0.74) & 1.07
(0.90-1.29) \\
& & Intellectual disability & 1.28 (0.85-1.91) & NR \\
& & Cerebral palsy & 1.41 (1.02-1.96) & 1.36 (0.78-2.38) \\
& & Epilepsy & 1.34 (1.05-1.71) & 1.21 (0.76-1.93) \\
& & Hearing loss & 1.45 (1.09-1.94) & 1.72 (0.97-3.02) \\
& & Arm fracture & 1.40 (1.19-1.65) & 3.27 (2.04-5.22) \\
UTI/pyelonephritis & 240,324 / 57,349 & Autism & 0.89
(0.85-0.92) & 0.98 (0.92-1.05) \\
& & ADHD & 0.74 (0.71-0.76) & 1.01 (0.95-1.06) \\
& & Speech-language disorder & 0.60 (0.59-0.61) & 0.98
(0.95-1.01) \\
& & Intellectual disability & 1.36 (1.27-1.47) & 1.13
(1.00-1.28) \\
& & Cerebral palsy & 1.19 (1.11-1.26) & 1.04 (0.93-1.15) \\
& & Epilepsy & 1.46 (1.39-1.53) & 1.19 (1.10-1.30) \\
& & Hearing loss & 0.91 (0.85-0.97) & 1.04 (0.94-1.16) \\
& & Arm fracture & 1.10 (1.07-1.14) & 1.13 (1.06-1.20) \\
RSV infection & 460,015 / 149,865 & Autism & 0.91
(0.89-0.94) & 0.93 (0.90-0.97) \\
& & ADHD & 0.85 (0.83-0.87) & 1.10 (1.07-1.14) \\
& & Speech-language disorder & 0.70 (0.69-0.71) & 1.02
(1.00-1.04) \\
& & Intellectual disability & 1.14 (1.07-1.22) & 0.93
(0.83-1.03) \\
& & Cerebral palsy & 0.96 (0.91-1.02) & 0.73 (0.67-0.80) \\
& & Epilepsy & 1.18 (1.14-1.23) & 0.94 (0.88-1.00) \\
& & Hearing loss & 0.71 (0.68-0.75) & 0.81 (0.75-0.87) \\
& & Arm fracture & 1.21 (1.18-1.24) & 1.11 (1.07-1.16) \\
Influenza & 314,636 / 127,676 & Autism & 0.90 (0.88-0.93) &
0.91 (0.87-0.95) \\
& & ADHD & 0.88 (0.86-0.90) & 1.05 (1.01-1.08) \\
& & Speech-language disorder & 0.75 (0.73-0.76) & 1.03
(1.01-1.05) \\
& & Intellectual disability & 0.90 (0.82-0.98) & 0.70
(0.61-0.79) \\
& & Cerebral palsy & 0.80 (0.74-0.87) & 0.52 (0.46-0.59) \\
& & Epilepsy & 1.17 (1.12-1.23) & 0.87 (0.81-0.93) \\
& & Hearing loss & 0.70 (0.66-0.75) & 0.79 (0.73-0.86) \\
& & Arm fracture & 1.20 (1.16-1.24) & 1.16 (1.11-1.21) \\
COVID-19 & 240,801 / 94,108 & Autism & 1.29 (1.24-1.33) &
1.22 (1.16-1.28) \\
& & ADHD & 1.37 (1.31-1.44) & 1.45 (1.36-1.54) \\
& & Speech-language disorder & 0.90 (0.89-0.92) & 1.20
(1.18-1.23) \\
& & Intellectual disability & 1.26 (1.11-1.44) & 0.82
(0.68-0.98) \\
& & Cerebral palsy & 0.89 (0.82-0.98) & 0.64 (0.56-0.72) \\
& & Epilepsy & 1.44 (1.36-1.53) & 1.04 (0.95-1.13) \\
& & Hearing loss & 0.75 (0.70-0.81) & 0.88 (0.80-0.97) \\
& & Arm fracture & 1.35 (1.29-1.42) & 1.20 (1.12-1.28) \\
Pneumonia & 331,838 / 101,444 & Autism & 0.80 (0.78-0.82) &
0.89 (0.84-0.93) \\
& & ADHD & 0.79 (0.77-0.81) & 1.05 (1.02-1.09) \\
& & Speech-language disorder & 0.65 (0.64-0.66) & 0.99
(0.97-1.01) \\
& & Intellectual disability & 1.44 (1.35-1.53) & 1.15
(1.04-1.26) \\
& & Cerebral palsy & 1.55 (1.47-1.63) & 1.23 (1.13-1.34) \\
& & Epilepsy & 1.50 (1.44-1.56) & 1.15 (1.08-1.23) \\
& & Hearing loss & 0.98 (0.93-1.03) & 1.06 (0.98-1.15) \\
& & Arm fracture & 1.11 (1.08-1.14) & 1.10 (1.05-1.15) \\
Pertussis & 13,094 / 3,033 & Autism & 0.83 (0.70-0.98) &
0.77 (0.58-1.00) \\
& & ADHD & 0.73 (0.64-0.82) & 1.05 (0.86-1.28) \\
& & Speech-language disorder & 0.62 (0.57-0.67) & 0.98
(0.87-1.11) \\
& & Intellectual disability & 0.89 (0.59-1.35) & 0.80
(0.42-1.53) \\
& & Cerebral palsy & 0.77 (0.54-1.10) & 0.51 (0.28-0.92) \\
& & Epilepsy & 1.42 (1.16-1.75) & 1.00 (0.69-1.45) \\
& & Hearing loss & 0.68 (0.51-0.92) & 0.90 (0.55-1.49) \\
& & Arm fracture & 1.12 (0.97-1.30) & 1.23 (0.94-1.61) \\
HSV infection & 38,830 / 12,702 & Autism & 0.76 (0.69-0.83)
& 0.82 (0.71-0.93) \\
& & ADHD & 0.82 (0.76-0.87) & 1.01 (0.92-1.11) \\
& & Speech-language disorder & 0.70 (0.67-0.73) & 1.04
(0.98-1.10) \\
& & Intellectual disability & 1.20 (0.98-1.46) & 0.86
(0.62-1.19) \\
& & Cerebral palsy & 1.70 (1.42-2.03) & 1.13 (0.84-1.52) \\
& & Epilepsy & 1.62 (1.44-1.84) & 1.03 (0.85-1.25) \\
& & Hearing loss & 0.90 (0.76-1.07) & 0.88 (0.68-1.14) \\
& & Arm fracture & 1.18 (1.09-1.28) & 1.10 (0.97-1.25) \\
Enterovirus infection & 125,491 / 39,922 & Autism & 0.95
(0.91-0.99) & 0.93 (0.86-0.99) \\
& & ADHD & 0.86 (0.83-0.90) & 1.04 (0.98-1.10) \\
& & Speech-language disorder & 0.77 (0.75-0.78) & 1.06
(1.02-1.09) \\
& & Intellectual disability & 1.24 (1.13-1.36) & 0.88
(0.75-1.02) \\
& & Cerebral palsy & 1.25 (1.16-1.35) & 1.09 (0.96-1.24) \\
& & Epilepsy & 1.62 (1.53-1.72) & 1.24 (1.12-1.37) \\
& & Hearing loss & 0.95 (0.88-1.03) & 1.00 (0.89-1.13) \\
& & Arm fracture & 1.13 (1.08-1.19) & 1.13 (1.05-1.22) \\
Viral gastroenteritis & 260,024 / 112,443 & Autism & 1.05
(1.02-1.09) & 1.05 (1.00-1.09) \\
& & ADHD & 0.99 (0.97-1.02) & 1.21 (1.17-1.25) \\
& & Speech-language disorder & 0.87 (0.86-0.89) & 1.16
(1.14-1.19) \\
& & Intellectual disability & 1.00 (0.93-1.09) & 0.83
(0.75-0.94) \\
& & Cerebral palsy & 0.86 (0.80-0.93) & 0.63 (0.57-0.70) \\
& & Epilepsy & 1.27 (1.21-1.33) & 0.97 (0.91-1.04) \\
& & Hearing loss & 0.97 (0.91-1.03) & 0.99 (0.92-1.08) \\
& & Arm fracture & 1.16 (1.13-1.20) & 1.15 (1.10-1.20) \\
Roseola & 38,795 / 12,265 & Autism & 0.74 (0.66-0.82) & 0.97
(0.83-1.13) \\
& & ADHD & 0.76 (0.70-0.82) & 0.97 (0.87-1.08) \\
& & Speech-language disorder & 0.56 (0.53-0.59) & 0.89
(0.83-0.96) \\
& & Intellectual disability & 0.81 (0.60-1.09) & 0.76
(0.48-1.20) \\
& & Cerebral palsy & 0.50 (0.36-0.70) & 0.60 (0.35-1.02) \\
& & Epilepsy & 0.92 (0.78-1.09) & 0.86 (0.67-1.11) \\
& & Hearing loss & 0.71 (0.57-0.89) & 1.03 (0.76-1.39) \\
& & Arm fracture & 1.17 (1.07-1.27) & 1.02 (0.89-1.17) \\
Varicella & 15,576 / 4,755 & Autism & 0.81 (0.69-0.95) &
0.91 (0.72-1.14) \\
& & ADHD & 0.78 (0.70-0.87) & 1.08 (0.93-1.26) \\
& & Speech-language disorder & 0.62 (0.57-0.67) & 0.96
(0.86-1.07) \\
& & Intellectual disability & 0.71 (0.48-1.07) & 1.14
(0.60-2.16) \\
& & Cerebral palsy & 0.44 (0.26-0.74) & NR \\
& & Epilepsy & 0.80 (0.63-1.01) & 0.50 (0.34-0.75) \\
& & Hearing loss & 0.98 (0.70-1.36) & 1.18 (0.72-1.92) \\
& & Arm fracture & 1.03 (0.91-1.16) & 1.09 (0.90-1.33) \\
Hand, foot, and mouth disease & 242,030 / 132,566 & Autism &
0.79 (0.77-0.82) & 0.77 (0.74-0.80) \\
& & ADHD & 1.00 (0.97-1.02) & 1.07 (1.04-1.10) \\
& & Speech-language disorder & 0.79 (0.77-0.80) & 0.96
(0.94-0.98) \\
& & Intellectual disability & 0.55 (0.49-0.61) & 0.42
(0.36-0.49) \\
& & Cerebral palsy & 0.39 (0.34-0.44) & 0.29 (0.25-0.34) \\
& & Epilepsy & 0.80 (0.75-0.85) & 0.63 (0.58-0.68) \\
& & Hearing loss & 0.76 (0.70-0.82) & 0.79 (0.72-0.87) \\
& & Arm fracture & 1.18 (1.15-1.22) & 1.10 (1.05-1.15) \\
Kawasaki disease & 10,796 / 1,825 & Autism & 0.79
(0.67-0.93) & 0.94 (0.68-1.29) \\
& & ADHD & 0.75 (0.65-0.87) & 0.73 (0.55-0.97) \\
& & Speech-language disorder & 0.51 (0.46-0.56) & 0.90
(0.76-1.06) \\
& & Intellectual disability & 0.95 (0.60-1.50) & 0.73
(0.34-1.54) \\
& & Cerebral palsy & 0.75 (0.46-1.20) & NR \\
& & Epilepsy & 0.76 (0.58-0.99) & 0.70 (0.39-1.24) \\
& & Hearing loss & 0.68 (0.48-0.96) & 0.95 (0.45-2.00) \\
& & Arm fracture & 0.95 (0.80-1.13) & 1.26 (0.85-1.85) \\
Diaper dermatitis (reference) & 602,639 / 409,413 & Autism &
1.06 (1.04-1.09) & 1.05 (1.03-1.08) \\
& & ADHD & 1.26 (1.24-1.28) & 1.36 (1.33-1.39) \\
& & Speech-language disorder & 1.05 (1.04-1.06) & 1.21
(1.19-1.22) \\
& & Intellectual disability & 0.85 (0.81-0.91) & 0.72
(0.67-0.77) \\
& & Cerebral palsy & 0.63 (0.60-0.67) & 0.50 (0.46-0.53) \\
& & Epilepsy & 1.02 (0.99-1.06) & 0.86 (0.83-0.90) \\
& & Hearing loss & 0.84 (0.81-0.88) & 0.88 (0.84-0.92) \\
& & Arm fracture & 1.09 (1.07-1.11) & 1.07 (1.04-1.10) \\
\end{longtable}

Unrestricted: exposure cohort did not require a well-child
visit (controls did). Bacterial meningitis, viral
meningitis/encephalitis, and HSV infection used 30
covariates (adding P92, K21, and Z82); others used the same
27 covariates as the primary analysis. CMV unrestricted
estimates are from a separate earlier analysis matched on 9
covariates (age at index, current age, sex, race, ethnicity,
P07, Q20-Q28, D80-D89, Q90-Q99). Unrestricted analyses were
not completed for febrile seizure, atopic dermatitis, and
iron-deficiency anemia.

\newpage

\hypertarget{etable-8.-e-values-for-statistically-significant-positive-associations}{%
\section*{eTable 8. E-Values for Statistically Significant
Positive
Associations}\label{etable-8.-e-values-for-statistically-significant-positive-associations}}

\begin{longtable}[]{@{}
  >{\raggedright\arraybackslash}p{(\columnwidth - 10\tabcolsep) * \real{0.3061}}
  >{\raggedright\arraybackslash}p{(\columnwidth - 10\tabcolsep) * \real{0.2755}}
  >{\raggedright\arraybackslash}p{(\columnwidth - 10\tabcolsep) * \real{0.1735}}
  >{\raggedleft\arraybackslash}p{(\columnwidth - 10\tabcolsep) * \real{0.0816}}
  >{\raggedleft\arraybackslash}p{(\columnwidth - 10\tabcolsep) * \real{0.0816}}
  >{\raggedleft\arraybackslash}p{(\columnwidth - 10\tabcolsep) * \real{0.0816}}@{}}
\toprule\noalign{}
\begin{minipage}[b]{\linewidth}\raggedright
Condition
\end{minipage} & \begin{minipage}[b]{\linewidth}\raggedright
Outcome
\end{minipage} & \begin{minipage}[b]{\linewidth}\raggedright
HR (95\% CI)
\end{minipage} & \begin{minipage}[b]{\linewidth}\raggedleft
FDR \emph{q}
\end{minipage} & \begin{minipage}[b]{\linewidth}\raggedleft
E-value, point
\end{minipage} & \begin{minipage}[b]{\linewidth}\raggedleft
E-value, CI limit
\end{minipage} \\
\midrule\noalign{}
\endhead
\bottomrule\noalign{}
\endlastfoot
Febrile seizure & Epilepsy & 8.33 (7.48-9.28) &
\textless.001 & 16.15 & 14.45 \\
CMV disease & Sensorineural hearing loss & 5.10 (3.67-7.09)
& \textless.001 & 9.67 & 6.79 \\
Bacterial meningitis & Epilepsy & 4.45 (3.33-5.93) &
\textless.001 & 8.36 & 6.12 \\
Bacterial meningitis & Cerebral palsy & 4.04 (2.91-5.61) &
\textless.001 & 7.55 & 5.27 \\
Bacterial meningitis & Sensorineural hearing loss & 3.91
(2.64-5.77) & \textless.001 & 7.28 & 4.73 \\
CMV disease & Cerebral palsy & 3.02 (2.09-4.38) &
\textless.001 & 5.49 & 3.59 \\
Viral meningitis/encephalitis & Epilepsy & 2.97 (1.92-4.59)
& \textless.001 & 5.38 & 3.24 \\
CMV disease & Intellectual disability & 2.68 (1.54-4.66) &
.001 & 4.80 & 2.45 \\
Viral meningitis/encephalitis & Cerebral palsy & 2.58
(1.42-4.66) & .004 & 4.59 & 2.20 \\
CMV disease & Epilepsy & 2.34 (1.67-3.30) & \textless.001 &
4.12 & 2.72 \\
Bacterial meningitis & Intellectual disability & 2.09
(1.21-3.62) & .02 & 3.60 & 1.71 \\
Sepsis (postneonatal) & Epilepsy & 2.07 (1.82-2.36) &
\textless.001 & 3.56 & 3.04 \\
Sepsis (postneonatal) & Cerebral palsy & 1.80 (1.57-2.06) &
\textless.001 & 3.00 & 2.51 \\
Sepsis (postneonatal) & Sensorineural hearing loss & 1.53
(1.32-1.78) & \textless.001 & 2.44 & 1.97 \\
COVID-19 & ADHD & 1.45 (1.36-1.54) & \textless.001 & 2.25 &
2.06 \\
CMV disease & Autism spectrum disorder & 1.45 (1.07-1.96) &
.03 & 2.25 & 1.35 \\
Bacterial meningitis & Speech-language disorder & 1.43
(1.22-1.68) & \textless.001 & 2.21 & 1.74 \\
Febrile seizure & Intellectual disability & 1.29 (1.06-1.56)
& .02 & 1.90 & 1.32 \\
CMV disease & Speech-language disorder & 1.25 (1.08-1.44) &
.005 & 1.80 & 1.38 \\
Enterovirus infection & Epilepsy & 1.24 (1.12-1.37) &
\textless.001 & 1.79 & 1.49 \\
Pneumonia & Cerebral palsy & 1.23 (1.13-1.34) &
\textless.001 & 1.77 & 1.52 \\
COVID-19 & Autism spectrum disorder & 1.22 (1.16-1.28) &
\textless.001 & 1.73 & 1.60 \\
Viral gastroenteritis & ADHD & 1.21 (1.17-1.25) &
\textless.001 & 1.70 & 1.61 \\
COVID-19 & Speech-language disorder & 1.20 (1.18-1.23) &
\textless.001 & 1.70 & 1.63 \\
UTI/pyelonephritis & Epilepsy & 1.19 (1.10-1.30) &
\textless.001 & 1.67 & 1.42 \\
Febrile seizure & Speech-language disorder & 1.19
(1.14-1.23) & \textless.001 & 1.66 & 1.54 \\
Febrile seizure & ADHD & 1.18 (1.11-1.25) & \textless.001 &
1.64 & 1.46 \\
Viral gastroenteritis & Speech-language disorder & 1.16
(1.14-1.19) & \textless.001 & 1.60 & 1.54 \\
Pneumonia & Epilepsy & 1.15 (1.08-1.23) & \textless.001 &
1.57 & 1.37 \\
Pneumonia & Intellectual disability & 1.15 (1.04-1.26) & .01
& 1.56 & 1.25 \\
Atopic dermatitis & Speech-language disorder & 1.15
(1.13-1.16) & \textless.001 & 1.56 & 1.52 \\
Sepsis (postneonatal) & Speech-language disorder & 1.12
(1.06-1.19) & \textless.001 & 1.49 & 1.31 \\
RSV infection & ADHD & 1.10 (1.07-1.14) & \textless.001 &
1.43 & 1.33 \\
Hand, foot, and mouth disease & ADHD & 1.07 (1.04-1.10) &
\textless.001 & 1.34 & 1.23 \\
Iron-deficiency anemia & Speech-language disorder & 1.06
(1.03-1.09) & \textless.001 & 1.32 & 1.23 \\
Atopic dermatitis & ADHD & 1.06 (1.04-1.08) & \textless.001
& 1.31 & 1.23 \\
Enterovirus infection & Speech-language disorder & 1.06
(1.02-1.09) & .001 & 1.31 & 1.19 \\
Pneumonia & ADHD & 1.05 (1.02-1.09) & .007 & 1.29 & 1.15 \\
Influenza & ADHD & 1.05 (1.01-1.08) & .02 & 1.27 & 1.12 \\
Atopic dermatitis & Autism spectrum disorder & 1.03
(1.01-1.06) & .03 & 1.21 & 1.09 \\
Influenza & Speech-language disorder & 1.03 (1.01-1.05) &
.01 & 1.20 & 1.10 \\
RSV infection & Speech-language disorder & 1.02 (1.00-1.04)
& .05 & 1.16 & 1.06 \\
\end{longtable}

Includes all condition-outcome pairs with a lower 95\% CI
limit above 1 and FDR \emph{q} \textless{} .05, ordered by
E-value.

\newpage

\hypertarget{etable-9.-complete-baseline-characteristics-for-the-cmv-comparison}{%
\section*{eTable 9. Complete Baseline Characteristics for the
CMV
Comparison}\label{etable-9.-complete-baseline-characteristics-for-the-cmv-comparison}}

\begin{longtable}[]{@{}
  >{\raggedright\arraybackslash}p{(\columnwidth - 12\tabcolsep) * \real{0.3514}}
  >{\raggedleft\arraybackslash}p{(\columnwidth - 12\tabcolsep) * \real{0.1081}}
  >{\raggedleft\arraybackslash}p{(\columnwidth - 12\tabcolsep) * \real{0.1622}}
  >{\raggedleft\arraybackslash}p{(\columnwidth - 12\tabcolsep) * \real{0.0721}}
  >{\raggedleft\arraybackslash}p{(\columnwidth - 12\tabcolsep) * \real{0.1081}}
  >{\raggedleft\arraybackslash}p{(\columnwidth - 12\tabcolsep) * \real{0.1261}}
  >{\raggedleft\arraybackslash}p{(\columnwidth - 12\tabcolsep) * \real{0.0721}}@{}}
\toprule\noalign{}
\begin{minipage}[b]{\linewidth}\raggedright
Characteristic
\end{minipage} & \begin{minipage}[b]{\linewidth}\raggedleft
CMV disease, before (n = 1,347)
\end{minipage} & \begin{minipage}[b]{\linewidth}\raggedleft
Controls, before (n = 3,090,897)
\end{minipage} & \begin{minipage}[b]{\linewidth}\raggedleft
SMD, before
\end{minipage} & \begin{minipage}[b]{\linewidth}\raggedleft
CMV disease, after (n = 1,345)
\end{minipage} & \begin{minipage}[b]{\linewidth}\raggedleft
Controls, after (n = 1,345)
\end{minipage} & \begin{minipage}[b]{\linewidth}\raggedleft
SMD, after
\end{minipage} \\
\midrule\noalign{}
\endhead
\bottomrule\noalign{}
\endlastfoot
Current age, mean (SD), y & 6.66 (4.47) & 7.45 (4.81) & 0.17
& 6.66 (4.48) & 6.65 (4.54) & 0.00 \\
Age at index, mean (SD), y & 0.18 (0.38) & 0.16 (0.36) &
0.06 & 0.18 (0.38) & 0.18 (0.39) & 0.02 \\
Male & 747 (55.5) & 1,581,797 (51.2) & 0.09 & 745 (55.4) &
765 (56.9) & 0.03 \\
Female \textsuperscript{a} & 597 (44.3) & 1,500,492 (48.5) &
0.08 & 597 (44.4) & 579 (43.0) & 0.03 \\
Hispanic or Latino & 367 (27.2) & 620,299 (20.1) & 0.17 &
366 (27.2) & 351 (26.1) & 0.03 \\
Not Hispanic or Latino & 786 (58.4) & 1,748,790 (56.6) &
0.04 & 785 (58.4) & 805 (59.9) & 0.03 \\
Unknown ethnicity \textsuperscript{a} & 194 (14.4) & 721,808
(23.4) & 0.23 & 194 (14.4) & 189 (14.1) & 0.01 \\
White & 564 (41.9) & 1,554,000 (50.3) & 0.17 & 563 (41.9) &
569 (42.3) & 0.01 \\
Black or African American & 343 (25.5) & 565,742 (18.3) &
0.17 & 343 (25.5) & 354 (26.3) & 0.02 \\
Asian & 102 (7.6) & 140,522 (4.5) & 0.13 & 102 (7.6) & 103
(7.7) & 0.00 \\
Native Hawaiian or Other Pacific Islander
\textsuperscript{a} & 31 (2.3) & 24,433 (0.8) & 0.12 & 31
(2.3) & ≤10 (0.7) & 0.13 \\
American Indian or Alaska Native \textsuperscript{a} & 11
(0.8) & 18,197 (0.6) & 0.03 & 11 (0.8) & ≤10 (0.7) & 0.01 \\
Other race \textsuperscript{a} & 112 (8.3) & 243,510 (7.9) &
0.02 & 112 (8.3) & 138 (10.3) & 0.07 \\
Unknown race \textsuperscript{a} & 184 (13.7) & 544,493
(17.6) & 0.11 & 183 (13.6) & 164 (12.2) & 0.04 \\
Short gestation/low birth weight (P07) & 364 (27.0) &
107,707 (3.5) & 0.69 & 362 (26.9) & 437 (32.5) & 0.12 \\
Slow fetal growth/fetal malnutrition (P05) & 115 (8.5) &
14,743 (0.5) & 0.40 & 113 (8.4) & 99 (7.4) & 0.04 \\
Newborn affected by maternal factors or delivery
complications (P00-P04) & 80 (5.9) & 56,212 (1.8) & 0.21 &
80 (5.9) & 102 (7.6) & 0.07 \\
Birth trauma (P10-P15) & ≤10 (0.7) & 9,473 (0.3) & 0.06 &
≤10 (0.7) & 15 (1.1) & 0.04 \\
Respiratory distress of newborn (P22) & 106 (7.9) & 15,036
(0.5) & 0.38 & 106 (7.9) & 105 (7.8) & 0.00 \\
Transitory endocrine/metabolic disorders of newborn
(P70-P74) & 108 (8.0) & 12,074 (0.4) & 0.39 & 108 (8.0) &
101 (7.5) & 0.02 \\
Other disturbances of cerebral status of newborn (P91) & 36
(2.7) & 6,218 (0.2) & 0.21 & 36 (2.7) & 33 (2.5) & 0.01 \\
Feeding problems of newborn (P92) \textsuperscript{a} & 163
(12.1) & 63,232 (2.0) & 0.40 & 163 (12.1) & 176 (13.1) &
0.03 \\
Congenital malformations of nervous system (Q00-Q07) & 173
(12.8) & 11,303 (0.4) & 0.52 & 171 (12.7) & 155 (11.5) &
0.04 \\
Congenital malformations of circulatory system (Q20-Q28) &
310 (23.0) & 64,870 (2.1) & 0.67 & 309 (23.0) & 343 (25.5) &
0.06 \\
Cleft lip and palate (Q35-Q37) \textsuperscript{a} & ≤10
(0.7) & 4,561 (0.1) & 0.09 & ≤10 (0.7) & 15 (1.1) & 0.04 \\
Congenital musculoskeletal malformations (Q65-Q79)
\textsuperscript{a} & 135 (10.0) & 105,946 (3.4) & 0.27 &
134 (10.0) & 146 (10.9) & 0.03 \\
Chromosomal abnormalities (Q90-Q99) & 32 (2.4) & 9,512 (0.3)
& 0.18 & 31 (2.3) & 36 (2.7) & 0.02 \\
Disorders of the immune mechanism (D80-D89) & 88 (6.5) &
1,958 (0.1) & 0.37 & 86 (6.4) & 49 (3.6) & 0.13 \\
Lack of expected normal physiological development (R62) &
301 (22.3) & 69,705 (2.3) & 0.64 & 300 (22.3) & 308 (22.9) &
0.01 \\
Symptoms concerning food and fluid intake (R63) & 304 (22.6)
& 93,265 (3.0) & 0.61 & 303 (22.5) & 300 (22.3) & 0.01 \\
Acute upper respiratory infection (J06) & 197 (14.6) &
199,340 (6.4) & 0.27 & 195 (14.5) & 202 (15.0) & 0.01 \\
Asthma (J45) & 24 (1.8) & 16,641 (0.5) & 0.12 & 24 (1.8) &
25 (1.9) & 0.01 \\
Nonsuppurative otitis media (H65) & 53 (3.9) & 28,354 (0.9)
& 0.20 & 53 (3.9) & 63 (4.7) & 0.04 \\
Suppurative/unspecified otitis media (H66) & 96 (7.1) &
88,933 (2.9) & 0.20 & 95 (7.1) & 100 (7.4) & 0.01 \\
Gastroesophageal reflux disease (K21) \textsuperscript{a} &
190 (14.1) & 156,747 (5.1) & 0.31 & 188 (14.0) & 221 (16.4)
& 0.07 \\
Family history of mental and behavioral disorders (Z81) &
≤10 (0.7) & 3,633 (0.1) & 0.10 & ≤10 (0.7) & ≤10 (0.7) &
0.00 \\
Family history of disabilities and chronic diseases (Z82)
\textsuperscript{a} & 45 (3.3) & 17,567 (0.6) & 0.20 & 45
(3.3) & 35 (2.6) & 0.04 \\
Well-child visit (any time to index) (Z00.1)
\textsuperscript{a} & 756 (56.1) & 3,090,897 (100.0) & 1.25
& 755 (56.1) & 1,345 (100.0) & 1.25 \\
Systemic antibacterials (ATC J01) & 441 (32.7) & 231,558
(7.5) & 0.66 & 440 (32.7) & 499 (37.1) & 0.09 \\
\end{longtable}

Abbreviations: CMV, cytomegalovirus; SMD, absolute
standardized mean difference.

Values are No.~(\%) unless otherwise indicated.
Before-matching denominators are infants with data in the
baseline window as reported by TriNetX (cohort sizes 1,386
and 3,228,691). TriNetX reports counts of 10 or fewer as 10;
these are shown as ≤10, and SMDs for them are approximate.
Race and ethnicity were as recorded in the electronic health
record.

\textsuperscript{a} Not included in the propensity score
model (shown to assess balance on variables not used for
matching). The well-child visit was required for all infants
before age 2 years; in the exposed group it could occur
after the CMV index date, so it appears in the covariate
window for only 56.1\%.

\newpage

\hypertarget{etable-10.-same-era-sensitivity-analysis-for-covid-19}{%
\section*{eTable 10. Same-Era Sensitivity Analysis for
COVID-19}\label{etable-10.-same-era-sensitivity-analysis-for-covid-19}}

\begin{longtable}[]{@{}
  >{\raggedright\arraybackslash}p{(\columnwidth - 10\tabcolsep) * \real{0.2883}}
  >{\raggedleft\arraybackslash}p{(\columnwidth - 10\tabcolsep) * \real{0.1261}}
  >{\raggedleft\arraybackslash}p{(\columnwidth - 10\tabcolsep) * \real{0.1261}}
  >{\raggedright\arraybackslash}p{(\columnwidth - 10\tabcolsep) * \real{0.1532}}
  >{\raggedright\arraybackslash}p{(\columnwidth - 10\tabcolsep) * \real{0.1532}}
  >{\raggedright\arraybackslash}p{(\columnwidth - 10\tabcolsep) * \real{0.1532}}@{}}
\toprule\noalign{}
\begin{minipage}[b]{\linewidth}\raggedright
Outcome
\end{minipage} & \begin{minipage}[b]{\linewidth}\raggedleft
COVID-19, events/at risk
\end{minipage} & \begin{minipage}[b]{\linewidth}\raggedleft
Controls, events/at risk
\end{minipage} & \begin{minipage}[b]{\linewidth}\raggedright
HR (95\% CI), COVID-19 vs same-era controls
\end{minipage} & \begin{minipage}[b]{\linewidth}\raggedright
HR (95\% CI), diaper dermatitis vs same-era controls
\end{minipage} & \begin{minipage}[b]{\linewidth}\raggedright
Calibrated HR (95\% CI)
\end{minipage} \\
\midrule\noalign{}
\endhead
\bottomrule\noalign{}
\endlastfoot
Autism spectrum disorder & 3,594/87,154 & 2,800/87,142 &
1.22 (1.17-1.29) & 1.14 (1.10-1.19) & 1.07 (1.01-1.14) \\
ADHD & 1,894/87,201 & 1,357/87,206 & 1.49 (1.39-1.60) & 1.53
(1.45-1.61) & 0.98 (0.90-1.07) \\
Speech-language disorder & 17,395/85,814 & 13,849/85,512 &
1.22 (1.19-1.25) & 1.25 (1.23-1.27) & 0.97 (0.95-1.00) \\
Intellectual disability & 179/87,178 & 179/87,167 & 1.01
(0.82-1.25) & 0.66 (0.55-0.79) & 1.54 (1.16-2.03) \\
Cerebral palsy & 355/87,117 & 495/87,082 & 0.67 (0.59-0.77)
& 0.50 (0.45-0.56) & 1.35 (1.13-1.60) \\
Epilepsy & 938/86,606 & 832/86,693 & 1.07 (0.98-1.18) & 0.82
(0.76-0.88) & 1.31 (1.16-1.48) \\
Sensorineural hearing loss & 667/86,823 & 735/86,846 & 0.86
(0.77-0.95) & 0.88 (0.82-0.95) & 0.97 (0.86-1.11) \\
Arm fracture (negative control) & 1,547/86,954 &
1,216/87,017 & 1.24 (1.15-1.34) & 1.08 (1.02-1.14) & 1.15
(1.05-1.27) \\
\end{longtable}

Both exposure and control cohorts were restricted to
children with a current age of 0-6 years (born from late
2019 onward). COVID-19 comparison: 93,329 exposed and
1,348,903 eligible controls before matching; 87,210 matched
pairs; median follow-up 1,001 vs 909 days. Diaper dermatitis
comparison: 225,532 and 1,216,700 before matching; 189,561
matched pairs; median follow-up 740 vs 688 days. Calibrated
HR = COVID-19 HR divided by the same-era diaper dermatitis
HR for the same outcome. Relative to its own
negative-control outcome (arm fracture), the COVID-19 autism
HR was 0.98 (95\% CI, 0.90-1.08) and the ADHD HR was 1.20
(95\% CI, 1.08-1.33). The proportional hazards test was
significant for autism (P \textless{} .001) and
speech-language disorder (P \textless{} .001) in the
COVID-19 comparison; HRs are averages over follow-up.
TriNetX analyses 6adfd51c027df79bdfadb92c and
6adfd5390a32ede36b4b2db0.

\end{landscape}

\hypertarget{efigure-1.-calibrated-hazard-ratios-for-autism-and-adhd}{%
\section*{eFigure 1. Calibrated Hazard Ratios for Autism and
ADHD}\label{efigure-1.-calibrated-hazard-ratios-for-autism-and-adhd}}

\includegraphics[width=6.5in,height=7.2in]{figs/eFigure1_calibrated_HR.png}

Ratio of each condition's HR to the diaper dermatitis HR for
the same outcome, with 95\% CIs (eMethods 6). Values above 1
indicate risk beyond that associated with a benign sick
visit.

\newpage

\hypertarget{efigure-2.-absolute-risk-differences-per-1000-children}{%
\section*{eFigure 2. Absolute Risk Differences per 1,000
Children}\label{efigure-2.-absolute-risk-differences-per-1000-children}}

\includegraphics[width=6.5in,height=7.2in]{figs/eFigure2_risk_difference.png}

Crude cumulative risk in exposed children minus matched
controls over 1 day to 9 years.

\newpage

\hypertarget{efigure-3.-primary-well-child-restricted-vs-unrestricted-design}{%
\section*{eFigure 3. Primary (Well-Child-Restricted) vs
Unrestricted
Design}\label{efigure-3.-primary-well-child-restricted-vs-unrestricted-design}}

\includegraphics[width=6.5in,height=7.2in]{figs/eFigure3_restricted_vs_unrestricted.png}

Each point is 1 condition; lines are 95\% CIs on each axis.
Points on the dashed diagonal have the same HR under both
designs.

\newpage

\hypertarget{efigure-4.-e-values}{%
\section*{eFigure 4. E-Values}\label{efigure-4.-e-values}}

\includegraphics[width=6.5in,height=7.2in]{figs/eFigure4_E_values.png}

E-values for statistically significant positive associations
(FDR \emph{q} \textless{} .05 and lower CI limit above 1).
Light bars, point estimate; dark bars, CI limit closest to
1.

\newpage

\hypertarget{ereferences}{%
\section*{eReferences}\label{ereferences}}

\begin{enumerate}
\def\labelenumi{\arabic{enumi}.}
\tightlist
\item
  Palchuk MB, London JW, Perez-Rey D, et al.~A global
  federated real-world data and analytics platform for
  research. \emph{JAMIA Open}. 2023;6(2):ooad035.
  doi:10.1093/jamiaopen/ooad035
\item
  Benjamini Y, Hochberg Y. Controlling the false discovery
  rate: a practical and powerful approach to multiple
  testing. \emph{J R Stat Soc Series B Stat Methodol}.
  1995;57(1):289-300. doi:10.1111/j.2517-6161.1995.tb02031.x
\item
  VanderWeele TJ, Ding P. Sensitivity analysis in
  observational research: introducing the E-value. \emph{Ann
  Intern Med}. 2017;167(4):268-274. doi:10.7326/M16-2607
\end{enumerate}

\end{document}
